\subsection{The specialization group and the unibranch case}

\begin{statement}{P4R.1}{Proved}
For a proper holomorphic degeneration of complex two-tori with
smooth total threefold and central fibre $mS$, where $S$ is
irreducible and reduced, the group $\pi_1(S)$ has an abelian
normal subgroup with cyclic quotient of order dividing $m$.
\end{statement}
\begin{proof}
We give the neighbourhood and group comparison explicitly.
The compatible triangulation of Lojasiewicz, Theorem~2,
\emph{Triangulation of semi-analytic sets}, p.~464 \cite{Lojasiewicz1964}, applies to
the smooth real analytic total space and its closed analytic
subset $S$. Its hypotheses hold because the total space is
countable and finite dimensional, and the collection consisting
of $S$ is locally finite and semianalytic. After barycentric
subdivision $S$ is a full subcomplex. In an open star of this
subcomplex, erase the barycentric coordinates at vertices
outside $S$ and normalize the remaining coordinates to sum
to one. Linear interpolation gives a strong deformation
retraction $r:V\to S$. The star can be chosen inside a
relatively compact neighbourhood of $S$.

Properness gives a small disc $\Delta_\epsilon$ with
$U=f^{-1}(\Delta_\epsilon)\subset V$: on the compact
complement of $V$ in a larger closed tube, $|f|$ has a
positive minimum. Thus the inclusion $i:S\to U$ and the
continuous restriction $r:U\to S$ satisfy $r\circ i=\id$.
This retraction comparison also follows from the statement
immediately after Definition~1 of Fernandez de Bobadilla--Kollar,
\emph{Homotopically trivial deformations}, p.~85 \cite{BobadillaKollar2012}; the preceding
construction supplies the particular comparison needed here.

Write $U^*=U\setminus S$. The proper smooth fibration
$U^*\to\Delta_\epsilon^*$ has fibre $T^4$. Its homotopy
sequence and $\pi_2(\Delta_\epsilon^*)=0$ give
\[
 \pi_1(U^*)=\Lambda\rtimes\langle t\rangle,
 \qquad \Lambda=\mathbb Z^4.
\]
The extension splits because its quotient is infinite cyclic.
The inclusion $U^*\to U$ induces a surjection on fundamental
groups. A loop can be perturbed off the real-codimension-two
smooth stratum of $S$ and off its higher-codimension singular
strata. The same general-position argument for a null-homotopy
disc avoids the singular strata and meets $S_{\rm reg}$ in
finitely many transverse points. Deleting little discs about
these intersections expresses the kernel as the normal closure
of the corresponding meridians. The smooth locus of the
irreducible $S$ is connected, so those meridians are conjugate,
with their positive orientation fixed by a complex transverse
coordinate.

At a smooth point of $S$, the defining map is a unit times
$z^m$. A positive meridian therefore maps to $m$ under
$\pi_1(U^*)\to\mathbb Z$. Let $A$ be the image of $\Lambda$
in $\pi_1(U)$. It is abelian and normal, and
$\pi_1(U)/A$ is cyclic with order dividing $m$.
The homomorphism $r_*:\pi_1(U)\to\pi_1(S)$ is surjective
because $r_*i_*=\id$. Its image $r_*(A)$ is abelian and
normal, and the induced quotient remains cyclic of order
dividing $m$. This proves the assertion for the actual support,
without replacing it by its normalization.
\end{proof}



\begin{statement}{P4R.2}{Proved}
In the setting of Part~4.2, every irreducible singular support
whose normalization is smooth and whose normalization morphism
is bijective has nonnegative topological Euler number.
Its Euler number is zero precisely when that normalization is
a minimal ruled surface over an elliptic curve or a Hopf surface.
\end{statement}
\begin{proof}
Let $\nu:T\to S$ denote this finite normalization morphism.
A continuous bijection from the compact $T$ to the Hausdorff
$S$ is a homeomorphism. It therefore induces an isomorphism
of fundamental groups and equality of Euler numbers. P4R.1
makes $\pi_1(T)$ virtually abelian.

The effective-conductor and canonical-torsion argument of
P4.12 applies to this same $T$ and proves $\kappa(T)=-\infty$.
The opened classification theorem used there leaves rational,
ruled, and class-VII minimal models, followed by point blowups.
Blowing up a smooth point preserves the fundamental group:
van Kampen applies to the common punctured four-ball with
simply connected boundary $S^3$. A ruled surface over a
genus-$g$ curve has the same fundamental group as that curve,
by the homotopy sequence of its $\mathbb P^1$-bundle.
For $g\ge2$ this group is not virtually abelian. Indeed a
finite-index subgroup corresponds to an unramified cover
of genus $1+d(g-1)>1$, and a genus-$h$ surface group maps
onto the nonabelian free group on $h$ generators by sending
the $b$ generators to the identity in its usual presentation.
Consequently $g\le1$.

A rational surface has positive Euler number. A ruled surface
over an elliptic curve, blown up $k$ times, has Euler number
$k$. A class-VII surface has Euler number $b_2\ge0$, and
each blowup increases it by one. These formulas give
$e(S)=e(T)\ge0$. Equality in the ruled case is exactly
$k=0$. Equality in class VII gives a minimal surface with
$b_2=0$. The conductor is nonempty because $S$ is singular
and nonnormal; P4.12 gives an effective conductor curve on
$T$. The precise class-VII classification and absence of
curves on Inoue surfaces verified in P4.8 therefore force
$T$ to be Hopf. The converse follows from the same Euler
formulas and the homeomorphism $\nu$.
\end{proof}

\subsection{The conductor pushout for ruled normalizations of higher genus}

\begin{statement}{P4R.3}{Proved}
In the setting of Part~4.2, let an irreducible singular fibre
support $S$ have smooth normalization $\nu:T\to S$ whose
minimal model is ruled over a curve of genus $g\ge2$.
Write $D$ for the conductor divisor on $T$ and
$C=\operatorname{Sing}(S)_{\rm red}$.
The inclusion $D_{\rm red}\to T$ induces an isomorphism on
fundamental groups, and the normalization pushout induces
$\pi_1(S)\simeq\pi_1(C)$. Moreover,
\[
 e(C)\ge0,\qquad e(S)>0.
\]
The assertion includes conductor components on different
ruling fibres that map onto the same curve in $C$.
\end{statement}
\begin{proof}
The canonical comparison of P4.12 gives
$K_T+D\equiv0$, with $D$ effective. We first determine the
reduced conductor as an embedded topological space; this
step does not require the ruling to descend to $S$.
Let $\pi:T\to T_0$ be the sequence of point blowdowns to
the minimal ruled surface $r_0:T_0\to B$, and put
$D_0=\pi_*D$. Then $K_{T_0}+D_0\equiv0$ and
$D_0\cdot F=2$ for a ruling fibre $F$.

There are three possibilities for the horizontal part of
$D_0$: a reduced bisection, two distinct reduced sections,
or twice one section. The first two are excluded directly.
For an irreducible horizontal component $H$ of coefficient
one, adjunction and the effectiveness of the other
components give
\[
 2p_a(H)-2=(K_{T_0}+H)\cdot H
       =-(D_0-H)\cdot H\le0.
\]
On the normalization of $H$, Riemann--Hurwitz for its
finite map of degree $d\ge1$ to $B$ gives
$2g(\widetilde H)-2\ge d(2g-2)>0$.
Since $p_a(H)\ge g(\widetilde H)$, these inequalities
are incompatible. Thus
\[
 D_0=2H+\sum_{j=1}^{v}a_jF_j,\qquad a_j\ge1,
\]
where $H$ is a section and the $F_j$ are distinct ruling
fibres. A horizontal curve of degree one is a section:
its normalization maps isomorphically to the smooth base,
and the resulting map to the projective bundle is a
section with closed image. In particular $H$ is smooth.

At one step $\beta:T'\to T''$ of the blowup sequence,
let $\mu$ be the multiplicity of the current conductor
$D''$ at the centre and $E$ the exceptional curve.
Write $D'=D''_{\rm strict}+aE$. The canonical blowup
formula gives
\[
 K_{T'}+D'
   =\beta^*(K_{T''}+D'')+(1+a-\mu)E.
\]
Intersection with $E$ and numerical triviality give
$a=\mu-1$. Effectiveness therefore places every centre
on the conductor. At a smooth conductor point of
coefficient one, the exceptional curve is absent from
the new conductor. At a smooth point of coefficient at
least two, it is a rational leaf attached at one point.
At an intersection of two conductor components it
subdivides their edge. Starting with $H$ and the
transverse fibres $F_j$, these operations preserve
normal crossings and make $D_{\rm red}$ a copy of $H$
with rational trees attached. No cycle is introduced.

The projection of $H$ to $B$ induces an isomorphism on
fundamental groups, and every attached rational tree
is simply connected. A point blowup leaves
$\pi_1(T)$ unchanged, as shown in P4R.2. The preceding
description also proves, at every blowup, that the
inclusion of the reduced conductor induces the same
isomorphism to the fundamental group of the surface.
Hence $\pi_1(D_{\rm red})\to\pi_1(T)$ is an
isomorphism.

We next verify the pushout for the actual singular
support. The morphism $\nu$ is finite, is an
isomorphism away from the conductor, and has smooth
source. Its nonisomorphism locus equals $C$: a normal
point has a neighbourhood identified with its
normalization and is consequently smooth. Conversely
every smooth point is normal. Thus the natural map
\[
 T\underset{D_{\rm red}}{\cup}C\longrightarrow S
\]
is a continuous bijection. Both source and target are
compact Hausdorff spaces, so it is a homeomorphism.
The compatible triangulation used in P4R.1 makes
$D_{\rm red}\subset T$ a subcomplex and hence a
cofibration. Replacing the map $D_{\rm red}\to C$ by
its mapping cylinder computes the same homotopy
pushout: the homotopy extension property for this
cofibration contracts the inserted cylinder without
altering the homotopy type. Van Kampen on that double
mapping cylinder therefore gives
\[
 \pi_1(S)=\pi_1(T)
       *_{\pi_1(D_{\rm red})}\pi_1(C)
       \simeq\pi_1(C).
\]
All three spaces in the amalgam are connected. This
argument retains every identification made by the
finite conductor map, including identifications
between horizontal and vertical branches.

P4R.1 makes this group virtually abelian. Let $c$ be
the number of irreducible components of $C$, let
$g_i$ be the genera of their smooth normalizations,
and let $b$ be the first Betti number of the incidence
graph of the finite point identifications. Replacing
each identification by a small tree joining its
preimages gives a homotopy model consisting of these
closed oriented surfaces joined along a connected
graph. Choosing a spanning tree gives
\[
 \pi_1(C)\simeq
   \Bigl(*_{i=1}^{c}\pi_1(C_i^{\rm norm})\Bigr)*F_b,
 \qquad
 e(C)=c+1-2\sum_i g_i-b.
\]
The formula also covers unibranch singular points:
their normalization is a local topological
homeomorphism and they contribute no graph edge.

A surface factor of genus at least two cannot occur
in a virtually abelian group, by the finite-cover
argument of P4R.2. Two nontrivial surface factors,
two free generators, or one elliptic factor together
with a free generator contain a nonabelian free
subgroup. Such a subgroup cannot lie in a virtually
abelian group: its intersection with a finite-index
abelian subgroup would have finite index and would
be abelian, whereas a finite-index subgroup of a
nonabelian free group has rank at least two. The
possibilities are therefore all $g_i=0$ with
$b\le1$, or exactly one $g_i=1$ with $b=0$.
The displayed Euler formula gives $e(C)\ge0$.
Additivity of compactly supported Euler numbers and
P4.13 now give
\[
 e(S)=e(S_{\rm reg})+e(C)\ge0.
\]

Finally equality is impossible here. It would force
$e(C)=0$, so the same formula gives $c=1$, $g_1=1$
and $b=0$. Every irreducible component of
$D_{\rm red}$ maps finitely and nonconstantly onto
this curve. A rational vertical component would
therefore give a nonconstant map from $\mathbb P^1$
to its genus-one normalization, contrary to
Riemann--Hurwitz. There are no vertical conductor
components. The initial conductor is consequently
$2H$, and there can have been no blowup: the first
centre would lie on $2H$ and would create a vertical
conductor component of coefficient one. Thus
$T=T_0$ and $D_{\rm red}=H$. It follows that
\[
 e(S_{\rm reg})=e(T\setminus H)
       =(4-4g)-(2-2g)=2-2g<0,
\]
contradicting P4.13. This proves strict positivity.
\end{proof}

P4R.3 proves strict positivity for its stated subclass.
The later classification P4.58 excludes every positive-Euler
ruled case, including this subclass, by the complete local
and global conductor arguments P4.54--P4.57 and P4R.18--P4R.20.
The earlier live status is retained in
\path{records/part4_root_superseded_status.md}.

\begin{statement}{P4R.4}{Proved}
In the setting of P4R.1, choose a positive small-loop
monodromy $A$ on $\Lambda=H_1(T^4;\mathbb Z)$.
For a suitable $\lambda\in\Lambda$, a positive
meridian of the reduced support is represented by
$\mu=\lambda t^m$ in $\Lambda\rtimes_A\langle t\rangle$.
The natural small-tube comparison gives
\[
 \pi_1(S)\simeq
 (\Lambda\rtimes_A\langle t\rangle)/\langle\!\langle\mu\rangle\!\rangle,
 \qquad
 \pi_1(S)/\operatorname{im}\Lambda\simeq\mathbb Z/m,
 \qquad
 b_1(S)=\dim_{\mathbb Q}\ker(A-I).
\]
For the support in P4R.3, a singular-locus incidence
graph that is a tree of rational normalization
components gives $e(S)\ge3$.
\end{statement}
\begin{proof}
We strengthen the neighbourhood comparison of P4R.1
without assuming that a preassigned tube is itself
one of the chosen open stars. Choose an open-star
neighbourhood $V_0$ retracting onto $S$, a small tube
$U_0\subset V_0$, a second open-star neighbourhood
$V_1$ with $S\subset V_1\subset U_0$, and a smaller
tube $U_1\subset V_1$. Properness supplies both tube
inclusions. The punctured-disc fibrations give an
isomorphism
\[
 \pi_1(U_1\setminus S)\longrightarrow
 \pi_1(U_0\setminus S):
\]
the fibres can be identified by radial parallel
transport and both base fundamental groups are
generated by the same positive loop. As in P4R.1,
the kernels of the maps to $\pi_1(U_i)$ are the
normal closures of the positive smooth-support
meridian. These meridians correspond under the
displayed isomorphism. Hence
$\pi_1(U_1)\to\pi_1(U_0)$ is an isomorphism.

This isomorphism factors through $\pi_1(V_1)$,
which is identified with $\pi_1(S)$ by its
retraction. Thus $\pi_1(S)\to\pi_1(U_0)$ is
surjective. The restriction to $U_0$ of the
retraction of $V_0$ proves injectivity. This
identifies $\pi_1(S)$ with the exact meridian
quotient, not just with a further quotient.

Projection to the cyclic factor sends $\mu$ to
$m$. After dividing by the image of $\Lambda$,
the remaining presentation is therefore exactly
$\langle t\mid t^m=1\rangle$. In the rational
abelianization the presentation becomes
\[
 \frac{(\Lambda/(A-I)\Lambda)\otimes\mathbb Q
       \ \oplus\ \mathbb Qt}
      {\mathbb Q(\overline\lambda+mt)}.
\]
The single displayed relation has nonzero
$t$ coefficient and reduces the dimension by
one. This proves the asserted formula for $b_1$.

In the rational-tree case of P4R.3 the same
pushout calculation gives $\pi_1(S)=1$.
The cyclic quotient then forces $m=1$, and the
$b_1$ formula gives $\det(I-A)\ne0$.
P4.13 gives
$e(S_{\rm reg})=\det(I-A)\ge1$.
A nonempty connected tree of $c\ge1$ rational
normalization components has Euler number
$c+1\ge2$, including any unibranch singular
points. Additivity gives $e(S)\ge3$.
\end{proof}

\subsection{Hopf normalizations and the conductor graph}

\begin{statement}{P4R.5}{Proved}
In the setting of Part~4.2, an irreducible singular
support whose normalization is smooth and has a Hopf
surface as its minimal model has nonnegative Euler
number. More precisely, for its reduced conductor
image $C$ one has $e(C)\ge0$. If the normalization is
obtained by $k$ point blowups and $v$ of the exceptional
curves belong to the conductor, then
\[
 e(S_{\rm reg})=k-v,\qquad
 e(S)=k-v+e(C)\ge0.
\]
\end{statement}
\begin{proof}
Let $T_0$ be the minimal Hopf surface. Every
irreducible curve $E$ on $T_0$ is smooth elliptic.
Indeed $b_2(T_0)=0$ gives $E^2=K_{T_0}\cdot E=0$,
so adjunction gives $p_a(E)=1$. A rational
normalization would define a nonconstant map
$\mathbb P^1\to T_0$. It lifts to the universal
cover $\mathbb C^2\setminus\{0\}$, where its two
coordinate functions are constant. This is
impossible. Thus the normalization genus is one
and the sum of the local delta invariants is zero.
Distinct curves are disjoint, since their
intersection number is zero and every local
intersection of distinct complex curves is
positive.

Push the effective conductor $D$ to $D_0$ on
$T_0$. The canonical comparison of P4.12 makes
$K_{T_0}+D_0$ torsion. The divisor $D_0$ is
nonempty: otherwise a positive power of $K_{T_0}$
would be trivial, contradicting the Kodaira
dimension $-\infty$ of a Hopf surface. Its
reduced support consists of $r\ge1$ disjoint
smooth elliptic curves. The blowup coefficient
$\mu-1$ calculation of P4R.3 applies without a
ruling. At a smooth conductor point it adds a
rational leaf precisely when $\mu\ge2$; at a
node it subdivides the corresponding edge.
Consequently $D_{\rm red}$ has $r$ connected
components, each a smooth elliptic root with a
rational tree attached. With $v$ rational
components it has exactly $v$ double points,
and hence $e(D_{\rm red})=v$. Since
$e(T)=k$, normalization away from the conductor
gives $e(S_{\rm reg})=k-v$. Each rational
component was created by one of the $k$ blowups,
so $k-v\ge0$.

Write $s$ for the number of connected components
of $C$, $c$ for its number of irreducible
components, and $g$ for the number of its
normalizations of genus one. Every normalized
component has genus zero or one: it receives a
finite nonconstant map from a normalized
conductor component. Each elliptic component
must receive an elliptic root, and distinct
image components require distinct roots.
Thus $c\ge g$ and $r\ge g$. Let $b$ be the sum
of the cycle ranks of the point-identification
graphs of the $s$ connected components. The
normalization formula gives
\[
 e(C)=c+s-2g-b.
\]

We retain the disconnected conductor in the
fundamental-group comparison. Its pushout has
one vertex for $T$ and one for each connected
component of $C$, with $r$ edges supplied by
the connected components of $D_{\rm red}$.
This graph is connected, because the finite
map $D_{\rm red}\to C$ is surjective.
The van Kampen presentation has $r-s$ stable
letters, in addition to the vertex groups.
The group of each conductor component is the
group of its elliptic root. Its map into $C$
factors, on that root, through the smooth
normalization of one image component; the
attached rational tree adds no fundamental
group generator. Kill the vertex group
$\pi_1(T)$ and the normal subgroups generated
by all normalized component groups of $C$.
Every conductor-group relation then disappears.
The resulting quotient is exactly
\[
 \pi_1(S)\twoheadrightarrow F_{b+r-s}.
\]
This is also the group of the homotopy-pushout
graph obtained by making the same quotients in
the normalization graph model of $C$. P4R.1
therefore gives $b+r-s\le1$. Substitution yields
\[
 e(C)\ge c+r-2g-1\ge-1.
\]

We exclude the only negative possibility by
the actual local hypersurface geometry.
The value $e(C)<0$ in these inequalities
forces $c=r=g$ and $b+r-s=1$. Every image
component then has elliptic normalization.
No rational conductor component can map
finitely onto one of them, by
Riemann--Hurwitz. Thus $D_{\rm red}$ consists
of disjoint smooth elliptic curves. This
conclusion permits blowups at coefficient-one
points; their exceptional curves are absent
from the conductor and do not affect the
following argument.

At a point $p\in C$, suppose there were two
distinct reduced curve germs $C_1,C_2$ of
$C$. Choose normalization preimages over $p$
whose conductor germs map onto $C_1,C_2$.
They must be distinct preimages, since the
reduced conductor has just one smooth germ
at each of its points. They correspond to
distinct irreducible analytic surface
branches $S_1,S_2$ of $S$ at $p$.
The support $S$ is a hypersurface in the
smooth total threefold. Its distinct
surface branches have a one-dimensional
intersection germ through $p$: two local
nonassociate irreducible equations in the
regular three-dimensional local ring have
height two. That intersection lies in the
conductor image of both branches. The
unique conductor germ on each chosen
normalization branch therefore has the
same image germ, giving $C_1=C_2$, a
contradiction.

Hence $C$ is locally unibranch and its
irreducible components are disjoint. Its
point-identification graphs have $b=0$
and $s=c=g$. The Euler formula gives
$e(C)=0$, contrary to the negative
possibility. We have proved $e(C)\ge0$,
and the asserted formulas finish the proof.
\end{proof}

\subsection{The integral pole divisor on a normalization resolution}

\begin{statement}{P4R.6}{Proved}
Let $S$ be an irreducible singular support in
Part~4.2, let $\nu:N\to S$ be its normalization,
and let $\rho:T\to N\to S$ use the minimal
resolution of $N$. There is an effective integral
divisor $D$ on $T$ such that
\[
 \omega_T(D)\simeq\rho^*\omega_S,
 \qquad \operatorname{Supp}D
       =\rho^{-1}(\operatorname{Sing}S)_{\rm red}.
\]
In particular $K_T+D$ is torsion, by P4.12.
The strict-transform conductor has its usual
positive integral coefficients. Every connected
exceptional configuration over a singular point
of $N$ has strictly positive coefficients in $D$.
Every such configuration meets the strict-transform
conductor.
\end{statement}
\begin{proof}
The reduced support is a Cartier hypersurface in
the smooth threefold, so $\omega_S$ is invertible.
Choose a local generator $\eta$ of this line
bundle. On $S_{\rm reg}$ it is an ordinary
holomorphic two-form; its pullback is a
meromorphic two-form on $T$. Its divisorial
orders are integers. Two generators differ by
a holomorphic unit on $S$, whose pullback is
still a unit. Consequently their negative
divisors define one integral divisor $D$ and
the displayed line-bundle isomorphism.
Finite duality, in the form
\[
 \nu_*\omega_N
  =\mathcal Hom_{\mathcal O_S}
       (\nu_*\mathcal O_N,\omega_S),
\]
identifies the generic strict-transform pole
ideal with the conductor. This is the same
normalization adjunction used in P4.12. Thus
the coefficients away from the exceptional
curves are the positive integral conductor
exponents; away from conductor and exceptional
curves the pullback has no zero or pole.
No assertion that $K_N$ alone is
$\mathbb Q$-Cartier is needed.

We prove positivity along an exceptional
configuration. Its intersection matrix $M$
is negative definite, and we give the local
comparison establishing the matrix property
used here. Choose a nonzero local function
on the normal surface $N$ vanishing at the
image point. Its pullback divisor is
$H+\sum_i n_iE_i$, with $n_i>0$ and $H$
effective without exceptional components.
The strict transform meets the exceptional
configuration, because a curve through the
image point has closure meeting its proper
inverse image. The divisor is principal
near the configuration, so
\[
 Mn=-h,\qquad h_i=H\cdot E_i\ge0,
 \qquad h\ne0.
\]
The off-diagonal entries of $M$ are
nonnegative. On a connected configuration,
writing a real vector as $v_i=n_iw_i$
gives the identity
\[
 v^tMv
 =-\sum_{i<j}M_{ij}n_in_j(w_i-w_j)^2
   -\sum_i n_ih_iw_i^2.
\]
It is strictly negative for every nonzero
$v$: equality forces all $w_i$ equal along
the connected graph and then forces that
common value to be zero. The same maximum
principle, applied to $M(n_iw_i)=-b$ with
$b_i\ge0$, shows that $-M^{-1}$ has
nonnegative entries, strictly positive on
each connected block. This also follows
directly by taking a negative minimum of
the $w_i$ in the displayed weighted
Laplacian identity.

Write $D=D_h+\sum_i d_iE_i$, with $D_h$
the strict conductor. Minimality of the
resolution gives
\[
 K_T\cdot E_i=2p_a(E_i)-2-E_i^2\ge0.
\]
Indeed the only exception for a negative
curve of arithmetic genus zero would be
a smooth rational exceptional curve of
self-intersection $-1$, which is absent
from a minimal resolution. Also
$D_h\cdot E_i\ge0$. Intersecting
$K_T+D=\rho^*K_S$ with the contracted
curves gives
\[
 Md=-b,
 \qquad b_i=(K_T+D_h)\cdot E_i\ge0.
\]
Hence $d_i\ge0$. A zero coefficient on
a connected block forces $b=0$ there.
Every curve in that block is then a
smooth rational $(-2)$-curve, and the
block is disjoint from the strict
conductor.

In fact, every exceptional block disjoint from
the strict conductor is forbidden, independently
of its coefficients. Finite duality identifies
the conductor on the normal $N$, up to the
invertible pullback of $\omega_S$, with the
rank-one reflexive sheaf $\omega_N$. Thus its
support has no isolated point. Any conductor
curve through the image point has strict
transform meeting the exceptional block:
properness makes the closure of its punctured
inverse image meet the fibre. Therefore the
image of a block disjoint from the strict
conductor is away from the conductor, and
$\nu$ is an isomorphism near it. The
corresponding point of $S$ is a normal isolated
hypersurface singularity. The local
Milnor-fibre argument of P4.12 applies
at this point itself. A local reduced
equation $g$ has an isolated critical
point and positive Milnor number. A
Morse perturbation supplies a vanishing
two-sphere of square $-2$ in its Milnor
fibre. The original map is a unit times
$g^m$; taking a local root of the unit,
each of its nearby local fibres is a
disjoint union of copies of a Milnor
fibre of $g$. Such a copy embeds into
the nearby four-torus. Every map
$S^2\to T^4$ is null-homotopic, so that
embedded sphere would have zero
homology class and square zero, a
contradiction. Thus every exceptional
block has positive pole coefficients.
This proves both effectiveness and the
asserted support equality.
\end{proof}

\begin{statement}{P4R.7}{Proved}
In Part~4.2, every irreducible singular
support whose normalization has a
resolution birationally ruled over a
curve of genus at least two has
nonnegative Euler number. The
normalization may be singular.
\end{statement}
\begin{proof}
Use the minimal resolution and effective
integral divisor of P4R.6. Push this
divisor to the minimal ruled model
$T_0$. Its intersection with a ruling
fibre is two. The adjunction and
Riemann--Hurwitz calculation in P4R.3
therefore gives exactly
\[
 D_0=2H+\sum_j a_jF_j,
 \qquad a_j\in\mathbb Z_{>0},
\]
with $H$ a smooth section. The reverse
blowup coefficient is the integral
number $\mu-1$. At a conductor crossing
one has $\mu\ge2$, so its exceptional
curve remains in the support and
subdivides the edge. At a smooth
conductor point the operation either
adds a rational leaf or leaves the
reduced curve unchanged. Thus
$A=\operatorname{Supp}D$ is connected,
is $H$ with rational trees attached,
and $\pi_1(A)\to\pi_1(T)$ is an
isomorphism.

Put $C=\operatorname{Sing}(S)_{\rm red}$.
The map $\rho$ is proper and is an
isomorphism outside $A$, and P4R.6
identifies $A$ with its full inverse
image of $C$. The continuous map
$T\cup_A C\to S$ is therefore a
compact-to-Hausdorff bijection and a
homeomorphism. The same cofibration
and van Kampen argument as P4R.3 gives
$\pi_1(S)\simeq\pi_1(C)$. This retains
the exceptional components that map
to points, as well as the strict
conductor components that map to
curves. P4R.1 and the normalization
graph calculation in P4R.3 give
$e(C)\ge0$. Finally P4.13 gives
$e(S_{\rm reg})\ge0$, so
$e(S)=e(S_{\rm reg})+e(C)\ge0$.
\end{proof}

\begin{statement}{P4R.8}{Proved}
In Part~4.2, every irreducible singular support
whose normalization has a resolution with Hopf
minimal model has nonnegative Euler number.
The normalization may be singular. For the
minimal normalization resolution, let $k$ be
the number of blowups from its Hopf minimal
model and $v$ the number of rational components
of the pole divisor. Then
\[
 e(S_{\rm reg})=k-v\ge0,\qquad
 e(\operatorname{Sing}S)\ge0.
\]
\end{statement}
\begin{proof}
Use $\rho:T\to N\to S$ and the effective
integral divisor $D$ of P4R.6. Put
$A=\operatorname{Supp}D$ and
$C=\operatorname{Sing}(S)_{\rm red}$.
The Hopf-curve and blowup calculations of
P4R.5 use only integrality, effectiveness,
and $K_T+D\equiv0$. They therefore apply
to this divisor: $A$ has $r\ge1$ connected
components, each a smooth elliptic root
with rational trees attached. Each of its
$v$ rational components came from one of
the $k$ blowups; $e(A)=v$, $e(T)=k$, and
$\rho$ is an isomorphism off $A$. This
proves the displayed formula for the
regular locus.

Every one-dimensional component of $C$
receives a nonconstant proper map from a
component of $A$, so its normalization
has genus zero or one. There are no
isolated points of $C$: an isolated
singular point of the hypersurface is
normal by the hypersurface $S_2$ property
and regularity in codimension one, and
P4R.6's vanishing-sphere argument excludes
it. Let $s,c,g,b$ respectively count the
connected components of $C$, its
irreducible components, its genus-one
normalizations, and the cycle ranks of
its point-identification graphs.
Then $e(C)=c+s-2g-b$, $r\ge g$, and
$c\ge g$.

The proper-map pushout $T\cup_A C=S$
and its cofibration justification are
exactly those of P4R.7. A root which
maps to a point gives the zero
homomorphism on its fundamental group;
a noncontracted root maps through the
normalization of one component of $C$.
Every rational tree is simply connected.
Consequently killing $\pi_1(T)$ and
the normalized component groups of $C$
in the disconnected van Kampen
presentation leaves the free quotient
$F_{b+r-s}$, as in P4R.5. Virtual
abelianness gives $b+r-s\le1$, hence
\[
 e(C)\ge c+r-2g-1\ge-1.
\]

A negative value forces $c=r=g$.
Every root is then noncontracted, maps
onto an elliptic normalization, and
supplies one of the $c$ image components.
Every rational component of $A$ is
contracted: a nonconstant map from its
normalization $\mathbb P^1$ to an
elliptic curve contradicts
Riemann--Hurwitz. Thus every exceptional
block is a rational subtree of one
rooted component of $A$. It meets that
root in exactly one point. At least one
intersection is forced by the final
assertion of P4R.6, since the roots are
the only strict conductor components;
two intersections would make a cycle
in the rooted tree. The exceptional
fibre is connected. One can see this
directly by Stein factorization of
$T\to N$: its finite birational factor
is an isomorphism because $N$ is normal.

The conductor on $N$ consequently has
one local analytic branch at each of
its points. At a singular point of $N$,
each local branch has a strict transform
meeting the exceptional fibre, whereas
there is exactly one such intersection
with the smooth root. Away from those
fibres the assertion follows from the
smooth disjoint roots themselves.

The local hypersurface argument of
P4R.5 now applies with $N$ in place of
its smooth resolution. Two different
curve germs of $C$ would require two
different normalization points, because
each point has just one conductor germ.
Their corresponding hypersurface
branches intersect in a curve germ.
That germ belongs to both conductor
images and forces the two chosen germs
to coincide. Thus $C$ is locally
unibranch with disjoint irreducible
components. This gives $b=0$, $s=c=g$,
and $e(C)=0$, a contradiction.
Therefore $e(C)\ge0$, proving the
assertion.
\end{proof}

\begin{statement}{P4R.9}{Proved}
In Part~4.2, every irreducible singular
support whose normalization has a
resolution birationally ruled over an
elliptic curve has nonnegative Euler
number. When the horizontal part of
the pole divisor on the minimal ruled
model is reduced, the normalization
is smooth and minimal and the support
has Euler number zero.
\end{statement}
\begin{proof}
Push the integral divisor of P4R.6 to
the minimal ruled surface $T_0\to B$.
Its horizontal degree is two. A reduced
horizontal component $H$ satisfies
\[
 0\le 2g(\widetilde H)-2
 \le 2p_a(H)-2
 =-(D_0-H)H\le0.
\]
The first inequality is
Riemann--Hurwitz over the elliptic
base. All inequalities are equalities:
every such component is smooth elliptic,
unramified over $B$, and disjoint from
the other components of $D_0$. Every
vertical fibre meets the horizontal
part, so no vertical component occurs.
The reduced-horizontal possibilities
are therefore one smooth elliptic
bisection or two disjoint elliptic
sections, each with coefficient one.

The reverse blowup formula $\mu-1$
then keeps $A=\operatorname{Supp}D$
a disjoint union of their smooth
strict transforms; each new exceptional
curve is absent from $A$. Every
exceptional curve of the normalization
resolution must belong to $A$, by
P4R.6. An exceptional component would
therefore be one of these disjoint
elliptic curves and would be disjoint
from the strict conductor, contrary
to P4R.6. Thus the normalization is
smooth. P4.22, with its hypotheses now
verified, proves minimality and
$e(S)=0$.

In the remaining case
$D_0=2H+\sum a_jF_j$ with $H$ a
smooth section. The same reverse
blowups make $A$ the elliptic root $H$
with rational trees attached. Its
inclusion in $T$ induces an
isomorphism on fundamental groups,
by the ruling and invariance under
point blowups. The resolution pushout
gives $\pi_1(S)=\pi_1(C)$ as in
P4R.7. The virtually abelian curve
group calculation gives $e(C)\ge0$,
and P4.13 gives $e(S_{\rm reg})\ge0$.
Their sum proves $e(S)\ge0$ also
in this remaining case.
\end{proof}
\subsection{The Hessian line and the tacnodal equality case}

\begin{statement}{P4R.10}{Proved}
Let $S$ be an irreducible reduced fibre support in Part~4.2.
Suppose that its reduced singular curve $C$ is irreducible and
unibranch at every point, that its normalization $j:E\to C\hookrightarrow Y$
is elliptic, and that its generic transverse germ is a tacnode
$x^2-y^4=0$. Let $C_{\mathrm{sch}}$ be the image conductor scheme.
Then
\[
 \chi(\mathcal O_{C_{\mathrm{sch}}})\leq0.
\]
Equality makes $C$ smooth and makes the transverse Hessian of
the reduced defining equation have rank one at every point of $C$.
\end{statement}
\begin{proof}
Let $P\subset\mathcal O_Y$ be the reduced ideal of $C$, and let
$I\subset\mathcal O_Y$ be the inverse image of the conductor ideal
of $S$. The conductor depth argument of P4.29 gives a
Cohen--Macaulay quotient $\mathcal O_Y/I$ of pure dimension one.
At the generic tacnode the normalization has two smooth branches;
its conductor is $(x,y^2)$ in the displayed coordinates.
Consequently $I$ contains $P^{(2)}$. Indeed a section of $P^{(2)}$
has zero image in $\mathcal O_Y/I$ at the generic point, and any
nonzero image would generate a finite-length submodule, contrary
to depth one. We obtain
\[
 0\longrightarrow J:=P/I\longrightarrow
 \mathcal O_{C_{\mathrm{sch}}}\longrightarrow\mathcal O_C
 \longrightarrow0.
\]
Here $J^2=0$, because $P^2\subset P^{(2)}\subset I$.
The module $J$ is therefore an $\mathcal O_C$-module. It is
torsionfree of rank one: a finite-support torsion submodule would
again be a finite-length submodule of $\mathcal O_Y/I$.

On $E$, put
\[
 K=\ker\bigl(j^*\Omega_Y\longrightarrow\Omega_E\bigr),
 \qquad\operatorname{im}(dj)=\Omega_E(-R).
\]
These are a rank-two bundle and a line bundle respectively.
The coefficient of $R$ at a point with local parametrization
$x_i=x_i(t)$ is $r=\min_i\operatorname{ord}_t x_i'(t)$.
We have $\deg K=\deg R$, since $\deg j^*\Omega_Y=0$
and $E$ is elliptic. Differentiation identifies
$G=P/P^{(2)}$ with a subsheaf of $j_*K$, as in P4.30.
Write $E_0\subset K$ for the image of its pullback to $E$.
This is a rank-two locally free subsheaf with finite-length
quotient.

The Hessian of a local reduced equation $f$ of $S$ annihilates
the tangent derivative of $C$, since all first derivatives of
$f$ vanish on $C$. It thus defines a quadratic tensor in
$\operatorname{Sym}^2K$. Under a change $f\mapsto u f$ its
restriction to $C$ multiplies by $u$, so globally it is a section
\[
 q\in H^0\bigl(E,\operatorname{Sym}^2K\otimes j^*\mathcal O_Y(S)\bigr).
\]
It has generic rank one. Let $L\subset K$ be the saturation of
its generic image line. On the smooth curve $E$ this is a line
subbundle. The tensor lies in $L^{\otimes2}\otimes
j^*\mathcal O_Y(S)$ everywhere, since that is a subbundle and
the assertion holds on a dense open set. Write its zero orders
as $d_p\geq0$. The normal line of $S$ has degree zero on $E$:
locally over the base disc the divisor $mS$ is principal.
It follows that
\[
 2\deg L=\sum_p d_p.
\]

Generically $I/P^{(2)}$ is exactly the line of the linear
conductor equation $x$, which is the Hessian line. Projection
to $K/L$ therefore gives an injection
\[
 J\longrightarrow j_*(K/L).
\]
To check the kernel at the exceptional points, a section whose
projection vanishes generically has torsion image in $J$; the
torsionfreeness just proved makes that image zero. Let $b_p$ be
the colength of the projected lattice $E_0\to K/L$ at $p$.
Let $q_p$ be the additional nonnegative length between $J$ and
the direct image of this normalized image lattice. Riemann--Roch
on $E$ now gives the exact formula
\[
 \chi(J)=\deg R-\frac12\sum_p d_p-\sum_p b_p-\sum_p q_p.
\]
We prove locally that $d_p/2+b_p\geq r_p$.

At a smooth point of $C$, $r_p=0$, so the claim follows from
nonnegativity. At a singular point put
$m=\min_i\operatorname{ord}_t x_i(t)$, so $r_p=m-1$.
First suppose the embedding dimension is three. Every defining
equation of $C$ has no linear term. Its gradient along the
parametrization has all entries of order at least $m$.
Every two-by-two minor of the conormal image consequently has
order at least $2m$. Since $K$ is saturated in the ambient free
rank-three module, a generator of $\bigwedge^2K$ has a unit
Pl\"ucker coordinate. Thus
\[
 \tau_p:=\operatorname{length}(K/E_0)_p\geq2m.
\]
Let $a_p=\operatorname{length}(L/(E_0\cap L))_p$.
The exact lattice sequence gives $\tau_p=a_p+b_p$.
Choose a primitive local generator $\ell$ of $L$ and write
$q=h\ell^2$ after trivializing the normal line, so
$\operatorname{ord}_t h=d_p$.
Every column of the Hessian is the differential of one of
the first derivatives of $f$. Those first derivatives belong
to $P$, so the columns lie in $E_0$. At least one ambient
coordinate of the primitive vector $\ell$ is a unit; its
corresponding column shows that $h\ell\in E_0$.
Therefore $d_p\geq a_p$, and
\[
 \frac{d_p}{2}+b_p\geq\frac{a_p}{2}+b_p
 \geq\frac{\tau_p}{2}\geq m>m-1=r_p.
\]

Now suppose the embedding dimension is two. Choose coordinates
with $P=(z,H(x,y))$, where $H$ is an irreducible reduced plane
equation with no linear term. Here $P^{(2)}=P^2$.
For an explicit verification, $z,H$ is a regular sequence, and
$P/P^2$ is free over the domain $\mathcal O_Y/P$.
If $s\notin P$ and $sv\in P^2$, reduction modulo $P$ gives
$v\in P$; reduction in this free conormal module then gives
$v\in P^2$. Thus $P^2$ is $P$-primary and equals its symbolic
square. The equation $f$ belongs to this square, so
\[
 f=A z^2+B zH+C H^2.
\]
Let $u$ be a primitive generator of the plane kernel of the
derivative row. We have $dH=w u$, with
$a=\operatorname{ord}_t w\geq m$, because both plane entries
of $dH$ have order at least $m$. In the basis $(dz,u)$ of $K$
we consequently have
\[
 E_0=\langle dz,w u\rangle,
 \qquad q=2A\,dz^2+2Bw\,dz\,u+2Cw^2u^2.
\]
Write the primitive generator of $L$ as
$\ell=\ell_z dz+\ell_u u$, and again $q=h\ell^2$ with
$\operatorname{ord}_t h=d_p$.
When $\ell_z$ is a unit, the projected conormal lattice has
colength $b_p=\min(a,\operatorname{ord}_t\ell_u)$.
The coefficient of $u^2$ in the displayed tensor gives
\[
 d_p+2\operatorname{ord}_t\ell_u\geq2a,
 \qquad d_p/2+b_p\geq a\geq m.
\]
The same inequality holds when $\ell_u=0$, since then $b_p=a$.
When $\ell_z$ is not a unit, $\ell_u$ is a unit. Projection of
$dz$ generates $K/L$, so $b_p=0$; the $u^2$ coefficient gives
$d_p\geq2a$, proving the same bound. These cases cover all
singular curve points and establish the claimed local inequality.

We conclude that $\chi(J)\leq0$. Since normalization gives
$\chi(\mathcal O_C)=-\sum\delta(C,p)$, the conductor exact
sequence yields
\[
 \chi(\mathcal O_{C_{\mathrm{sch}}})
 =-\sum_p\delta(C,p)+\chi(J)\leq0.
\]
Equality forces every delta invariant to vanish, so $C=E$ is
smooth. Then $R=0$, $E_0=K$, and all normalization descent
defects vanish. The exact formula reduces to
$\chi(J)=-\tfrac12\sum d_p$, so equality also forces
$d_p=0$ everywhere. The Hessian has rank one at every point.
The program \texttt{code/part4/tacnode\_hessian.py} verifies the
two lattice inequalities for the stated bounds, the exact
rank-one Hessian on the singular parametrization
$(t^2,t^3,0)$, and the simple Hessian zero of the permitted
local model $zx^2+xy^2$. Its JSON and text outputs retain
the bounds; the preceding inequalities prove the unrestricted
statement.
\end{proof}

\begin{statement}{P4R.11}{Proved}
Let $S$ be an irreducible singular support in Part~4.2 whose
normalization has ruled resolution, and suppose $e(S)=0$.
In the double-horizontal-divisor case with generic transverse
tacnode, the normalization is a smooth minimal elliptic ruled
surface. The conductor image is a smooth elliptic curve, its
reduced preimage is an unramified covering of degree two,
and every local germ is a product tacnode.
\end{statement}
\begin{proof}
Use $T\to N\to S$, the horizontal section $H$, and
$C=\operatorname{Sing}(S)_{\mathrm{red}}$ as in P4R.9.
The fundamental-group and Euler argument of P4.29 gives an
irreducible curve $C$ with elliptic normalization, unibranch
at every point. All rational curves in the resolved conductor
are contracted, since none maps nonconstantly to an elliptic
curve. Thus $H$ is the only strict conductor component.
Its finite map to $C$ has degree two, and its genus $g$ is at
least one by Riemann--Hurwitz.

Every exceptional curve of $T\to N$ is vertical for the ruling
$T\to H$. The normalized-graph construction of P4.29 descends
this ruling to $N\to H$. Its Leray calculation proves
$R^1\rho_{N*}\mathcal O_T=0$, and hence
$\chi(\mathcal O_N)=1-g$. Finite duality identifies the
conductor on $N$ with $\omega_N\otimes\nu^*\omega_S^{-1}$.
The latter line twist is torsion by P4.12. Serre duality,
the rational-singularity comparison and Riemann--Roch on $T$
therefore give $\chi(I)=1-g$ and
$\chi(\mathcal O_{D_{\mathrm{sch}}})=0$.
The exact conductor square, with $\chi(\mathcal O_S)=0$,
then gives
\[
 \chi(\mathcal O_{C_{\mathrm{sch}}})=g-1.
\]
P4R.10 makes this at most zero. Thus $g=1$, $C$ is smooth,
and the reduced defining equation has rank-one Hessian along
all of $C$. Riemann--Hurwitz makes $H\to C$ unramified.

We verify the local analytic conclusion. Over any point of
$C$, its two preimages on $H$ cannot be identified in $N$:
each exceptional block is vertical for the descended ruling
and hence meets only one point of the section $H$.
There are therefore two local irreducible branches of $S$.
The holomorphic Morse lemma with parameters writes its reduced
equation as $X^2+g(y,z)$. Its two analytic factors are monic
linear polynomials in $X$, by Weierstrass factorization, and
their coefficients sum to zero. Each branch is a graph and
is smooth. Their intersection has support exactly the smooth
curve $C$. Consequently, in coordinates with $C=\{X=y=0\}$,
the factor coefficient is a unit times a power of $y$.
The generic tacnode fixes that power to two. Extracting a square
root of the unit in the $y$ coordinate gives
$X^2-y^4=0$ times the $z$ disc. This also verifies the actual
normalization and conductor morphisms locally.
The normalization was already smooth away from the conductor
by P4R.6. It is therefore smooth everywhere, and P4.22 makes
this elliptic ruled surface minimal.
\end{proof}

\subsection{The components outside the pole divisor in a ruled fibre}

\begin{statement}{P4R.12}{Proved}
Let $T\to N\to S$ be the minimal normalization resolution of an
irreducible support in Part~4.2 with ruled minimal model of positive
base genus. Let $r:T\to B_0$ be its ruling, and let $D$ be the
effective integral pole divisor of P4R.6. Every component $E$ of a
reducible ruling fibre which is absent from $\operatorname{Supp}D$
is a smooth rational curve of self-intersection $-1$. It meets
exactly one pole component, transversely once, and that component
has pole coefficient one. In the double-horizontal-section case,
$E$ is disjoint from the horizontal section. Distinct such
components are disjoint; each is a leaf of the reduced ruling fibre.
These assertions retain the multiplicities of the ruling fibre
and do not require the multiplicity of $E$ to be one.
\end{statement}
\begin{proof}
A reduced horizontal pole divisor occurs here only over an
elliptic base, by P4R.3, and P4R.9 then makes the normalization
smooth and minimal. In that case there is no reducible ruling
fibre. We may therefore work in the double-horizontal-section
case throughout the proof.
A ruling fibre on $T$ is obtained from a smooth $\mathbb P^1$
by the actual sequence of point blowups. Its reduced support is
a normal-crossing tree of smooth rational curves. Each component
of a reducible fibre has negative self-intersection. One can
verify the latter assertion without a contraction criterion:
if $F_b=\sum m_iE_i$ is the full fibre, then
$0=F_bE_i=m_iE_i^2+\sum_{j\ne i}m_jE_jE_i$;
connectedness and reducibility make the last sum positive.

For $E$ outside the pole support, effectiveness and adjunction give
\[
 0\le DE=-K_TE=2+E^2.
\]
Thus $E^2$ is $-1$ or $-2$. In the second case $DE=0$,
so $E$ is disjoint from the full pole support. The map to $S$
then identifies a neighbourhood of this compact curve with a
neighbourhood in $S_{\rm reg}$, by P4R.6. The local cyclic
base change which removes the multiplicity of the support is
unramified over this regular locus. Its restriction to $E$ is
a disjoint union of copies of $\mathbb P^1$, because
$\pi_1(\mathbb P^1)=0$. A chosen copy has the same normal
degree $-2$. A smooth local trivialization over a sufficiently
small parameter disc transports it into a nearby four-torus,
still with self-intersection $-2$. Every map from $S^2$ to that
torus is null-homotopic, since its universal cover is
$\mathbb R^4$; its homology class and self-intersection must
therefore be zero. This contradiction leaves $E^2=-1$ and
$DE=1$.

All pole coefficients and all intersection multiplicities are
positive integers at an intersection. Hence $DE=1$ says exactly
that there is one intersection, transverse and on a coefficient-one
component. A horizontal pole section of coefficient two cannot
be that component.

It remains to exclude an intersection of two components $E,E'$
outside the pole divisor. Both have square $-1$, and both meet
the pole support. In the full ruling fibre, let their positive
multiplicities be $m,m'$. The two fibre-intersection equations
would give
\[
 m=m'+a,\qquad m'=m+b,
\]
where $a,b$ are positive sums of the multiplicities of their
other vertical neighbours; their respective pole neighbours make both
sums positive. Adding these equations gives $a+b=0$, a
contradiction. Thus no two outside components meet. Each has
exactly the one pole neighbour already found and is a leaf.
\end{proof}

\subsection{Restriction to an elliptic conductor component}

\begin{statement}{P4R.13}{Proved}
In Part~4.2, let an irreducible singular support $S$ have a
normalization with ruled resolution. If its reduced image
conductor has a component with elliptic normalization, then
the normalization of $S$ is a smooth minimal elliptic ruled
surface, the image conductor is one smooth elliptic curve,
and $e(S)=0$. Every local germ along this curve is the product
of its generic transverse germ with a disc. In particular no
rational conductor component is attached to that curve.
\end{statement}
\begin{proof}
Let $\rho_N:T\to N$ be the minimal normalization resolution,
let $\nu:N\to S$ be the finite normalization, and put
$\rho=\nu\rho_N$. The reduced-horizontal case is already
proved by P4R.9. We therefore use the actual double-horizontal
pole divisor $D$ of P4R.6. Its support is a smooth section $H$
with rational trees attached. The pushout and curve-group
calculation of P4R.7--P4R.9 show that the reduced image
conductor $C$ is connected, that its point-identification
graph has no cycles, and that it has exactly one normalized
component of positive genus, of genus one. Denote the
corresponding reduced component by $C_E$ and its normalization
by $j:E\to C_E\hookrightarrow Y$. In particular every point
of $C_E$ is unibranch. Every other normalized component of
$C$ is rational.

A rational component of $D$ cannot map nonconstantly to $E$.
Thus $H$ maps nonconstantly to $E$. Write $g=g(H)\ge1$.
At a general point of $C_E$, every normalization branch has
conductor exponent two, because the coefficient of $H$ in
$D$ is two. The intrinsic conductor exponent of a singular
plane branch is at least two, and each additional branch
increases it by its positive intersection multiplicity.
The possibilities are therefore a single ordinary cusp,
two smooth branches with intersection multiplicity two, or
three smooth branches with pairwise intersection multiplicity
one. These are $A_2,A_3,D_4$, respectively, and $H\to E$
has degree $1,2,3$ in the three cases.

We first construct a map which detects global functions on
the entire conductor, including its rational components.
Since $H$ is not contracted, P4.33 proves that $N$ has
rational singularities. Put $L_N=\nu^*\omega_S$.
Finite duality identifies its conductor ideal with
\[
 I_N=\omega_N\otimes L_N^{-1}.
\]
The canonical trace is an isomorphism
$\rho_{N*}\omega_T=\omega_N$ for rational analytic
singularities. We use Greb, Remark~2.1, p.~415
\cite{Greb2011}. The hypotheses hold here: $N$ is normal
and Cohen--Macaulay, $\rho_N$ is a proper resolution,
P4.33 gives $R^1\rho_{N*}\mathcal O_T=0$, and its
one-dimensional fibres give vanishing of the higher
structure-sheaf direct images. The integral pole identity
of P4R.6 and projection formula consequently give
\[
 \rho_{N*}\mathcal O_T(-D)
 =\rho_{N*}(\omega_T\otimes\rho_N^*L_N^{-1})
 =I_N.
\]
These are equalities of ideal subsheaves of $\mathcal O_N$:
the pole identity uses the pullback of the same residue
section which defines finite duality, and the trace agrees
with it on the birational isomorphism locus. Equality there
determines the maps between these torsionfree sheaves.
The divisor ideal sequence therefore gives an injection
\[
 \mathcal O_N/I_N\longrightarrow\rho_{N*}\mathcal O_D.
\]
The finite conductor square gives an injection from the
image conductor scheme. Composing them yields the actual
morphism of coherent $\mathcal O_S$-algebras
\[
 \mathcal O_{C_{\rm sch}}\longrightarrow
 \nu_*(\mathcal O_N/I_N)\longrightarrow\rho_*\mathcal O_D,
 \qquad\hbox{both arrows injective}.                         \tag{P4R.13.1}
\]

We show that a global function on $D$ is determined by
its value at the generic doubled horizontal section.
At every point blowup $\beta:T'\to T''$ in the actual
map to the minimal ruled surface, the pole divisors obey
$D'=\beta^*D''-E_\beta$. The elementary direct images are
\[
 \beta_*\mathcal O(E_\beta)=\mathcal O,
 \quad R^1\beta_*\mathcal O(E_\beta)=0,
 \quad\beta_*\mathcal O=\mathcal O,
 \quad R^1\beta_*\mathcal O=0.
\]
For the first equality, a meromorphic function allowed
a pole only on $E_\beta$ descends off the centre and
extends across that point by Hartogs. For the vanishing,
the two blowup charts, or their exceptional-curve
filtration, have successive terms
$\mathcal O_{\mathbb P^1}(j)$ with $j\ge-1$;
their first cohomology vanishes. Equivalently one applies
$0\to\mathcal O\to\mathcal O(E_\beta)
\to\mathcal O_{E_\beta}(-1)\to0$ to the usual
blowup structure-sheaf calculation. Projection formula,
the divisor ideal sequence and the absence of second
direct images then give
\[
 \beta_*\mathcal O_{D'}=\mathcal O_{D''},
 \qquad R^1\beta_*\mathcal O_{D'}=0.                         \tag{P4R.13.2}
\]
These particular maps identify $H^0(D,\mathcal O_D)$
with $H^0(D_0,\mathcal O_{D_0})$, where
\[
 D_0=2H_0+V_0,
 \qquad V_0=\sum_j a_jF_j
\]
on the minimal ruled surface. The kernel of restriction
from $D_0$ to $2H_0$ is $\mathcal O_{V_0}(-2H_0)$.
Its successive layers on each thick fibre $a_jF_j$ are
$\mathcal O_{\mathbb P^1}(-2)$, since $F_j^2=0$ and
$H_0F_j=1$. They have no global sections. Restriction
to $2H_0$ is therefore injective on global functions.
This double section is Cohen--Macaulay and has no
zero-dimensional associated point, so a section which
vanishes generically on it vanishes everywhere.

Let $C_{E,\rm sch}$ be the primary component of
$C_{\rm sch}$ supported on $C_E$: locally its ideal
is obtained by localizing the conductor ideal at the
generic point of $C_E$ and contracting. Primary
decomposition of the pure one-dimensional conductor
quotient makes these ideals coherent and compatible
on overlaps. Restriction is a surjection
\[
 \mathcal O_{C_{\rm sch}}\longrightarrow
 \mathcal O_{C_{E,\rm sch}}.
\]
Its map on $H^0$ is injective. Indeed, a global element
of its kernel maps by (P4R.13.1) to a function on $D$
which vanishes generically on $2H$; (P4R.13.2) and the
preceding paragraph make it zero. Its map on $H^1$
is surjective, because the kernel is supported in
dimension at most one and has no second cohomology.
We have proved
\[
 \chi(\mathcal O_{C_{\rm sch}})
 \le\chi(\mathcal O_{C_{E,\rm sch}}).                       \tag{P4R.13.3}
\]

We evaluate the right side using the actual primary
ideal, without removing the rational components from
the left side. In type $A_2$ the primary conductor
is reduced: a pure one-dimensional quotient which
is generically reduced has no nilpotent submodule
of finite support. Thus its Euler characteristic
is $-\sum_p\delta(C_E,p)\le0$. In type $D_4$ the
primary ideal is the symbolic square of the ideal
of $C_E$, by exactly the generic-conductor and
depth-one argument of P4.31. P4.30 gives the same
nonpositive bound, with equality only for smooth
$C_E$.

In type $A_3$, the primary ideal contains that
symbolic square and its nilpotent module is
torsionfree of rank one on $C_E$. The proof of
P4R.10 applies to this primary quotient: its
differentiation map, saturated Hessian line and
projection to the quotient conormal line depend
only on this ideal and on the actual equation of
$S$. The presence of other conductor components
changes none of those maps. More explicitly its
Euler characteristic is
\[
 -\sum_p\delta(C_E,p)+\deg R
 -\frac12\sum_p d_p-\sum_p b_p-\sum_p q_p\le0,
\]
with the same ramification divisor $R$, Hessian
zero orders $d_p$, projected lattice colengths
$b_p$ and descent defects $q_p$ as in P4R.10.
The local inequalities there prove this bound
at every singular point of $C_E$. Equality
makes $C_E$ smooth and its Hessian nowhere zero.

On the other hand P4.33 gives
$\chi(\mathcal O_{C_{\rm sch}})=g-1$. Thus
\[
 0\le g-1
 =\chi(\mathcal O_{C_{\rm sch}})
 \le\chi(\mathcal O_{C_{E,\rm sch}})\le0.                    \tag{P4R.13.4}
\]
All these quantities vanish. We obtain $g=1$,
$C_E=E$ smooth, and an unramified map $H\to E$.
In type $A_3$ its transverse Hessian has rank
one everywhere on $E$.

It remains to exclude attachments to this smooth
elliptic component. Put
\[
 K=\mathcal I_E/\mathcal I_E^2,
 \qquad N_E=\mathcal O_Y(S)|_E.
\]
Both $\deg K$ and $\deg N_E$ are zero. The first
follows from the conormal exact sequence,
$c_1(TY)E=0$, and $\deg\Omega_E=0$; the second
follows because a positive multiple of $S$ is
principal in a tube of this fibre.

In type $A_2$ let $L\subset K$ be the saturated
quadratic-image line and put $Q=K/L$. The nonzero
quadratic tensor is a section of
$L^2\otimes N_E$, so $2\deg L\ge0$.
The cubic jet on the null direction is an
intrinsic section of $Q^3\otimes N_E$.
Here is the coordinate check, including points
where the quadratic tensor vanishes. A change
of a second-order lift changes the cubic
coefficient by the polarization of the
quadratic tensor. That polarization vanishes
on the null direction. A change of tangential
lift contributes a tangential derivative of
the rank-one quadratic tensor; its evaluation
twice on the null direction is again zero.
Changing a reduced defining equation multiplies
the tensor by its transition unit. Thus the
stated line bundle and its holomorphic section
are defined on all of $E$. The section is
nonzero generically by type $A_2$, so
$-3\deg L\ge0$. Both degrees are zero,
and both sections are nowhere zero. The
parametrized Morse lemma and a cube root
of the residual unit now give the product
equation $X^2+Y^3=0$ at every point of $E$.

In type $A_3$ (P4R.13.4) already makes the
quadratic tensor nowhere zero, whence
$\deg L=\deg Q=0$. Its cubic null jet is
identically zero. Removing the quadratic
direction gives an intrinsic residual fourth
jet in $Q^4\otimes N_E$. In local coordinates
its relevant terms and coefficient are
\[
 a x^2+bxy^2+cy^4,
 \qquad q_4=c-\frac{b^2}{4a},\qquad a\ne0.                 \tag{P4R.13.5}
\]
Substitution of the critical coordinate
$x=-by^2/(2a)+O(y^3)$ proves the formula.
The change $x\mapsto x+vy^2$ sends $b$ to
$b+2av$ and $c$ to $c+bv+av^2$, leaving
$q_4$ unchanged. Changes of normal frames
and of equation give the claimed
$Q^4\otimes N_E$ transition. Tangential
lift changes do not affect a function
whose residual order is four. This section
is nonzero generically and has degree zero;
it never vanishes. The residual function
is $y^4$ times a unit, giving a product
tacnode at every point of $E$.

In type $D_4$, the equation belongs to
$\mathcal I_E^3$. This ideal is primary
because $E$ is smooth, so generic order
three proves membership everywhere.
The cubic normal jet belongs to
$\operatorname{Sym}^3K\otimes N_E$.
Its discriminant is a section of
$(\det K)^6\otimes N_E^4$, of degree
zero. It is nonzero generically, hence
nowhere zero. The three distinct normal
tangent directions give three smooth
holomorphic branches locally along $E$.
After straightening the first two to
$x=0,y=0$, the third is $ax+by=0$
with both coefficients units. Setting
$X=ax,Y=-by$ gives the product equation
$XY(X-Y)=0$.

In all three cases the entire singular
locus near every point of $E$ is $E$
itself. A rational conductor tree could
not meet it. Connectedness of $C$ then
forces $C=E$. The displayed local
product forms make $N$ smooth along
its conductor, and P4R.6 already makes
it smooth elsewhere. Hence $T=N$.
A finite normalization cannot contract
a curve, and a rational curve cannot
map nonconstantly to $E$. Thus $D=2H$
has no vertical component. The first
blowup of $2H_0$ in a nonempty actual
blowup sequence would create a vertical
pole component of coefficient one;
later blowups never remove its strict
transform from the pole support.
Consequently no blowup occurs and $N$
is minimal elliptic ruled. Finally
\[
 e(S)=e(N)-e(H)+e(E)=0.
\]
The exact coefficient changes, discriminant
weights and degree identities are executed
in \texttt{code/part4/elliptic\_primary\_restriction.py};
its JSON and text outputs are retained.
\end{proof}

\subsection{The smooth critical curve of transverse type \texorpdfstring{$A_3$}{A3}}

\begin{statement}{P4R.14}{Proved}
Let $S$ be an actual reduced support in Part~4.2. At a point
where the entire reduced critical locus of its local defining
function is one smooth curve, with generic transverse type
$A_3$, holomorphic coordinates put that function in one of
the forms
\[
 X^2+Y^4,\qquad X(Y^2+XZ).                              \tag{P4R.14.1}
\]
The normalization germs are respectively two smooth branches,
or a smooth plane and an $A_1$ quadric. No positive-genus
exceptional curve occurs over such a point.
\end{statement}
\begin{proof}
Choose a coordinate $z$ on the critical curve and normal
coordinates $x,y$. Its quadratic transverse Hessian has generic
rank one. Diagonalization of the symmetric matrix over
$\mathbb C\{z\}$, as in P4.28, gives the quadratic term
$h(z)x^2$. The cubic coefficient on the null direction
vanishes identically, by the generic transverse type. With
weights two for $x$ and one for $y$, the weight-four terms are
\[
 h(z)x^2+a(z)xy^2+b(z)y^4,\qquad
 \Delta(z)=4h(z)b(z)-a(z)^2.                             \tag{P4R.14.2}
\]
Every remaining term has larger weight and $\Delta$ is
nonzero generically.

We first prove that $\Delta$ is a unit. A zero persists,
counted with multiplicity, under small independent constant
perturbations of $h,a,b$, by the argument principle on a
small parameter disc. Generic such perturbations make the
zeros simple and disjoint from $h=0$. Indeed $h=\Delta=0$
also requires $a=0$, which the independent parameters avoid;
where $h\ne0$, the derivative of $\Delta$ in the $b$
parameter is $4h\ne0$. These perturbations preserve the
critical axis and its transverse $A_3$ germ near the fixed
boundary. At a simple zero with $h\ne0$, the parametrized
Morse lemma gives
\[
 X^2+y^4U(y,z),\qquad U(0,0)=0,\quad U_z(0,0)\ne0.
\]
The actual coordinate $Z=U(y,z)$ gives $X^2+y^4Z$.
P4.16 supplies a compact spherical class with nonzero square
in its smoothing. The boundary comparison of P4.28 applies
here: near the boundary critical curve, the parametrized
Morse lemma and a fourth root of a unit give $X^2+Y^4$;
the tangential boundary coordinate makes the nonzero levels
submersive. Proper boundary transport therefore carries the
sphere to the actual nearby torus. Its universal cover is
$\mathbb R^4$, so a sphere has zero homology class and zero
square, a contradiction. Thus $\Delta$ is a unit.

For $h(0)\ne0$, Morse splitting and a fourth root of the
residual unit give the first form in (P4R.14.1). Suppose
$h(0)=0$. Then $a(0)\ne0$. The holomorphic function
\[
 \frac{f(y^2v,y,z)}{y^4}
\]
restricts at $y=0$ to $h(z)v^2+a(z)v+b(z)$. Its finite
root $-b(0)/a(0)$ at the origin is simple. The holomorphic
implicit-function theorem gives a root $v=v(y,z)$, and
division by the monic equation of its graph gives
\[
 X=x-y^2v(y,z),\qquad f=Xg.
\]
On $X=0$, the function $g$ is $y^2$ times a unit. Its
parameter Morse centre in $y$ is therefore zero there and
its critical value is zero. The parametrized Morse lemma,
with parameters $X,z$, gives the exact form
\[
 f=X\bigl(Y^2+XH(X,z)\bigr).                            \tag{P4R.14.3}
\]
The remaining critical value is divisible by $X$ by its
vanishing on $X=0$. The function $H(0,z)$ is not identically
zero, since the generic transverse Hessian has rank one;
at the point under consideration $H(0,0)=0$.

Put $G=2H+XH_X$. Off $X=0$, the critical equations of
(P4R.14.3) are $Y=0,H_z=0,G=0$. We prove $H_z(0,0)\ne0$.
Otherwise $H_z$ is a nonzero plane germ not divisible by
$X$, because $H(0,z)$ has finite order at least two.
No component of $H_z=0$ is contained in $G=0$: it would
give an additional critical-curve branch off $X=0$, whereas
the entire critical locus is the stated smooth axis. Hence
$(H_z,G)$ is an isolated complete intersection of positive
intersection multiplicity.

Replace $H$ by $H-\eta z-\varepsilon$. Its plane critical
equations are
\[
 H_z-\eta=0,\qquad G-2\eta z-2\varepsilon=0.              \tag{P4R.14.4}
\]
Their solutions are the fibres of the holomorphic map
\[
 \Psi(X,z)=\left(H_z,\frac{G-2zH_z}{2}\right).
\]
Its zero fibre has the positive local degree just proved.
Choose a small compact representative with no boundary zero.
That degree is unchanged on a sufficiently small open
parameter ball. The universal equations in (P4R.14.4) have
parameter Jacobian determinant two, so regular parameter
values give a nonempty finite set of transverse solutions.
We also choose such a value outside the images under $\Psi$
of $X=0$ and $H_X=0$. These exceptional images have real
four-dimensional measure zero: each source curve has
countably many smooth real two-dimensional strata, and the
smooth image of each has measure zero. The germ $H_X$ is
not identically zero. Indeed in that case $H$ would depend
only on $z$ and both $H_z$ and $G=2H$ would vanish on the
whole line $z=0$, contrary to the isolated intersection.
The two exceptional images cannot exhaust the open ball
on which the positive degree is retained. This argument
uses no analyticity assertion for the image of a nonproper
map.

At each chosen solution, the critical value is
\[
 f_{\eta,\varepsilon}=-\frac{X^3H_X}{2}\ne0.              \tag{P4R.14.5}
\]
For the indicated order of the two equations, put
\[
 J=\det\frac{\partial(G-2\eta z-2\varepsilon,\ H_z-\eta)}
                  {\partial(X,z)}.
\]
The $Y$-Hessian entry is $2X$, and the remaining Hessian
block at a critical point is
\[
 \begin{pmatrix}
 X(3H_X+XH_{XX})&X^2H_{Xz}\\
 X^2H_{Xz}&X^2H_{zz}
 \end{pmatrix}.
\]
Its determinant is $X^3J$, and the full Hessian determinant
is $2X^4J\ne0$. Thus these are isolated holomorphic Morse
points with nonzero critical values. Near the fixed boundary,
$H(0,z)$ remains nonzero under the small perturbation, so the
critical axis retains its transverse $A_3$ type. The same
proper boundary comparison transports a Morse vanishing
sphere, of square $-2$, into an actual nearby torus. This is
impossible. We have proved $H_z(0,0)\ne0$.

The holomorphic coordinate $Z=H(X,z)$, keeping $X,Y$, now
gives the second form of (P4R.14.1). Its normalization has
the plane $X=0$ and the normal quadric $Y^2+XZ=0$. The
quadric's minimal resolution has one rational $(-2)$-curve.
The first form factors as $(X+iY^2)(X-iY^2)$ and its two
normalization branches are smooth. This proves the stated
normalization conclusions.

The distinction between a pole coefficient and this
torus-smoothing conclusion can be checked by an exact
retained local test. Take
\[
 f=x\bigl[y^2+xz(x-z)(x-2z)\bigr].
\]
It has only the smooth $A_3$ critical axis $x=y=0$. The
second normalization branch is the double plane branched
over four distinct lines. Blowing up their common point
and normalizing gives a double cover of the exceptional
$\mathbb P^1$ branched in four points, hence an elliptic
curve of square $-2$. Its discrepancy is $-1$ and the
order of $x$ on it is one, giving pole coefficient two
for the displayed total reduced hypersurface. It is excluded
here by the polar deformation, not by the coefficient alone.
For $H=z(x-z)(x-2z)$, the deformation $H\mapsto H-1$ has
the exact Groebner basis
\[
 3x-5z^4-14z,\qquad 25z^6+50z^3-2.
\]
It gives six simple solutions; adjoining $x$, $H_x$ or
$J$ to the critical ideal gives the unit ideal in each
case. For $H\mapsto H-\varepsilon$ all six scale by
$\varepsilon^{1/3}$ and converge to the origin. The
identities and these exact rational calculations are
executed in \texttt{code/part4/smooth\_a3\_polar.py},
with JSON and text output retained.
\end{proof}

\subsection{A nonzero Hessian and the entire critical curve}

\begin{statement}{P4R.15}{Proved}
At every critical-curve point of an actual reduced support in
Part~4.2 with nonzero ambient Hessian, holomorphic coordinates
give the exact product germ
\[
                       X^2+Y^m,\qquad m\ge2.
\]
In particular its entire reduced critical curve is smooth.
At a singular point of the reduced critical curve, the ambient
Hessian vanishes. This assertion allows different generic
transverse types on the different branches before the exclusion.
\end{statement}
\begin{proof}
A nonzero complex symmetric Hessian has a vector on which its
quadratic form is nonzero. The parametrized holomorphic Morse
lemma in that direction gives the exact equation
\[
                        f=X^2+g(Y,Z).
\]
The plane function $g$ is nonzero because the support is
reduced. Factor it in the convergent power-series ring as
\[
             g=u\prod_{i=1}^{r}h_i^{m_i},
             \qquad u(0)\ne0,\quad m_i\ge1,
\]
with distinct reduced irreducible factors $h_i$ vanishing at
the point. At least one exponent is at least two: a reduced
plane function has an isolated critical point, while the
critical locus here contains a curve.

We exclude $r\ge2$ by an actual deformation of these factors.
For two distinct factors $h_i,h_j$, the holomorphic map
\[
 (h_i,h_j):(\mathbb C^2,0)\longrightarrow(\mathbb C^2,0)
\]
has isolated zero fibre and positive finite local degree.
On a sufficiently small compact representative that degree
is retained for small target values. Regular such values
give a nonempty transverse intersection of
$h_i-s_i=0$ and $h_j-s_j=0$. Independently choose the
other constants $s_k$ so that at least one of these
intersection points avoids every other zero factor. This
is possible because the chosen intersection set is finite
and each forbidden $s_k$ belongs to a finite set of values.
All parameters can be arbitrarily small.

At the selected intersection, the two transverse factors
are holomorphic coordinates and every other factor is a
unit. Taking an actual local $m_i$th root of their unit
product and multiplying the first coordinate by it gives
the germ
\[
                        X^2+Y^{m_i}Z^{m_j}.
\]
P4.18 excludes every such germ by its compact spherical
class of nonzero square. The comparison here is to the
original torus degeneration, as follows. Choose a small
Milnor ball before moving the factors. Near its boundary,
the distinct reduced factors are smooth and disjoint.
The moved factor with exponent $m_i\ge2$ retains exactly
the transverse product type $A_{m_i-1}$ there; a factor
of exponent one remains a regular point of the surface.
Parametrized Morse coordinates and the appropriate root
of the remaining unit identify these boundary germs.
Solving the tangential boundary coordinate makes the
nonzero levels submersive. The proper boundary transport
used in P4.28 therefore identifies the needed nearby
levels and transports the spherical class into the actual
nearby torus. Its square cannot be nonzero there. Hence
there is only one factor.

Absorb the unit by writing $h=u^{1/m}h_1$, so that
$g=h^m$. The reduced plane function $h$ is irreducible.
For a singular $h$, its gradient has an isolated zero,
and the finite positive local degree of
$(h_Y,h_Z)$ gives nonempty finite regular fibres under
small constant changes of that gradient. Equivalently,
subtracting a sufficiently small generic linear form
from $h$ gives isolated Morse critical points. Add a
small constant to make their critical values nonzero.
Call the resulting function $h_s$, and choose one such
critical point with value $v\ne0$ and nonsingular Hessian
$Q$. For $f_s=X^2+h_s^m$, the critical value and Hessian
are exactly
\[
 f_s=v^m\ne0,\qquad
 \operatorname{Hess}(f_s)
   =\begin{pmatrix}2&0\\0&m v^{m-1}Q\end{pmatrix}.
\]
Thus the full Hessian determinant is
$2m^2v^{2m-2}\det Q\ne0$. The moving curve $h_s=0$
retains its transverse $A_{m-1}$ type on the fixed
boundary. The same proper comparison transports the
Morse vanishing sphere, of square $-2$, to a nearby
torus, a contradiction. Consequently $h$ is smooth,
and using it as the coordinate $Y$ gives the asserted
product equation.

The Hessian identity is executed symbolically for the
arbitrary positive integer $m$ in
\texttt{code/part4/hessian\_product.py}. Its retained
output also checks the moving-node example
$X^2+(YZ-\varepsilon)^m$ at a nonzero critical value.
\end{proof}


\subsection{The tacnode and transverse-line critical curve}

\begin{statement}{P4R.16}{Proved}
In the remaining positive ruled case with uniform transverse
type $A_3$, the three-branch point of the critical curve cannot
be the union of a plane tacnode and an independent transverse
line. Its local conductor defect is at least eleven, whereas
the total defect in the contracted-section case is nine.
\end{statement}
\begin{proof}
The analytic ring calculation in P4.40 supplies actual ambient
coordinates for this curve. Its reduced ideal is
\[
 P=(xz,yz,h),\qquad h=y^2-x^4.
\]
Its three parametrizations are $(t,t^2,0)$, $(t,-t^2,0)$,
and $(0,0,t)$. The two tangent branches have intersection
multiplicity two, and their union meets the independent line
with intersection length one. The normalization sequence
therefore gives $\delta=3$. Every branch is immersive, so
its ramification divisor has coefficient zero.

We first determine the symbolic square and the actual
equation. The plane-tacnode ideal $(z,h)$ is a reduced
complete intersection. Its square is unmixed and agrees
with the square of each prime at its generic point; hence
its symbolic square is its ordinary square. Intersecting
with the square of the line ideal gives
\[
 P^{(2)}=(z,h)^2\cap(x,y)^2=P^2+(zh).
\]
Indeed an element of $(z,h)^2$ is $Az^2+Bzh+Ch^2$.
Its last two terms belong to $(x,y)^2$. The first belongs
to that ideal precisely when $A$ does, because $z$ is a
nonzerodivisor modulo $(x,y)^2$. The resulting three
generators $x^2z^2,xyz^2,y^2z^2$ are already in $P^2$.
Furthermore
\[
 xzh=(xz)h,\qquad yzh=(yz)h,\qquad
 z^2h=(yz)^2-x^2(xz)^2.
\]
The class of $zh$ modulo $P^2$ is therefore annihilated
by the ambient maximal ideal. It is nonzero, since it has
a nonzero cubic term and $P^2$ has order at least four.
Thus $P^{(2)}/P^2$ is exactly one copy of the residue field.

Let $f$ be the reduced defining function of the support.
Its generic transverse multiplicity is two, so
$f\in P^{(2)}$ and
\[
                       f=\kappa zh+f_0,\quad f_0\in P^2,
                       \quad\kappa\in\mathbb C.
\]
At $(t,t^2,0)$ with $t\ne0$, use $z,h$ as transverse
coordinates. The quadratic transverse part of $f$ has the
form
\[
 t^2a(t)z^2+(\kappa+t b(t))zh+c(t)h^2,
\]
where $a,b,c$ are holomorphic at $t=0$. The transverse
Hessian has rank one, so its determinant vanishes for
all sufficiently small nonzero $t$. Its constant term at
zero is $-\kappa^2$ (up to the harmless common convention
for the quadratic-form matrix). Hence $\kappa=0$, and
we have proved the exact inclusion $f\in P^2$.

Let $K_i$ be the saturated rank-two kernel of the derivative
of the $i$th parametrization, and let $E_i\subset K_i$ be
the span of the differentiated generators of $P$. On the
two tangent branches choose the bases
\[
                dz,\quad dy\mp2t\,dx.
\]
The generator matrices have columns
\[
 \begin{pmatrix}t&\pm t^2&0\\0&0&\pm2t^2\end{pmatrix}.
\]
Their elementary-divisor orders are $(1,2)$. On the
independent line use $dx,dy$; the matrix is
\[
                   \begin{pmatrix}t&0&0\\0&t&0\end{pmatrix},
\]
with orders $(1,1)$.

For completeness, the rank-one quadratic estimate of
P4.42 applies to these exact lattices as follows. Write
$E=t^a\mathbb C\{t\}e_1\oplus
t^b\mathbb C\{t\}e_2$, and write the saturated Hessian
line as $\ell=u e_1+v e_2$, with $(u,v)$ primitive.
Since $f\in P^2$, the tensor $t^d\ell^2$ belongs to
$\operatorname{Sym}^2E$. Therefore
\[
 \frac d2\ge\max(a-\operatorname{ord}u,
                         b-\operatorname{ord}v).
\]
Projection of $E$ to $K/\ell$ has colength
$b_{\rm proj}=\min(a+\operatorname{ord}v,
b+\operatorname{ord}u)$. Adding these two expressions
gives $d/2+b_{\rm proj}\ge a+b$; the identity follows
by subtracting the same two candidate quantities, with
the zero-component cases read directly. The local defect
formula of P4.42 consequently gives
\[
 \delta+\sum_i(d_i/2+b_i-r_i)+q
       \ge3+(1+2)+(1+2)+(1+1)=11.
\]
Here $q\ge0$ is the further descent colength, retained
rather than discarded from the definition. P4.44 contracts
the horizontal section, and P4.33 and (4.42a) give total
defect nine. The displayed lower bound contradicts that
total. This excludes the stated analytic curve germ.
\end{proof}

The symbolic-square identities, all three normalized
conormal matrices, and the rank-one determinant obstruction
are checked exactly in
\texttt{code/part4/a3\_tacnode\_line.py}.


\subsection{The nonordinary three-branch tacnode sector}

\begin{statement}{P4R.17}{Proved}
In the remaining positive ruled case with uniform generic transverse
type $A_3$, every nonordinary three-branch critical-curve point has
local conductor defect strictly greater than nine. In particular,
the delta-four nonplanar cases have the explicit lower bounds
$13$, $21/2$, $11$, and $23/2$ below. The only three-branch germ
compatible with the total defect budget nine is the ordinary
nonplanar union of three coordinate axes.
\end{statement}
\begin{proof}
Use the notation of P4R.10 and P4.42: \[
 D(C,f)=\delta(C)+
        \sum_i\left(\frac{d_i}{2}+b_i-r_i\right)+q,\qquad q\ge0.
\] Here \(d_i\) is the order of the saturated rank-one Hessian, \(b_i\) is the projected conormal colength, and \(r_i=m_i-1\) for a branch of multiplicity \(m_i\).

\par\medskip\noindent\textbf{1. Bounds without an ordinary-square assumption}\par\label{bounds-without-an-ordinary-square-assumption}

Let \(E_i\subset K_i\) be the normalized conormal lattice, with elementary divisors \((t^{e_1},t^{e_2})\), \(e_1\le e_2\). For the saturated Hessian line \(L_i\), put \[
 a_i=\operatorname{length}\bigl(L_i/(L_i\cap E_i)\bigr).
\] The exact lattice sequence gives \(a_i+b_i=e_1+e_2\). Every first derivative of \(f\) belongs to the radical ideal \(P\) of the critical curve. Thus every Hessian column belongs to \(E_i\). Write the Hessian as \(ut^{d_i}\ell_i^2\), with \(u\) a unit and \(\ell_i\) primitive. One coordinate of \(\ell_i\) is a unit; its column gives \(t^{d_i}\ell_i\in E_i\), so \(d_i\ge a_i\). Also \(b_i\ge e_1\), since \(E_i\subset t^{e_1}K_i\). Consequently \[
                  \frac{d_i}{2}+b_i
                    \ge e_1+\frac{e_2}{2}.           \tag{A}
\] This remains valid when a symbolic-square term occurs in \(f\).

Since the curve is nonplanar, \(P\) has no linear equation. Every defining gradient therefore has order at least \(m_i\) on branch \(i\), and \(b_i\ge m_i\). P4R.15 gives zero ambient Hessian, so \(f\) has order at least three and \(d_i\ge m_i\). For total multiplicity \(m=\sum m_i\), \[
                  D(C,f)\ge\delta(C)+3+\frac m2.       \tag{B}
\] Thus \(m\ge5\) gives \(D\ge19/2>9\). Only multiplicities three and four need inspection.

\par\medskip\noindent\textbf{2. Three smooth branches: the complete excess-two classification}\par\label{three-smooth-branches-the-complete-excess-two-classification}

Put \[
 B=\bigoplus_{i=1}^3\mathbb C\{t_i\},\quad
 F^j=\bigoplus_{i=1}^3t_i^j\mathbb C\{t_i\},\quad
 A_0=\mathbb C_{\rm diag}+F^1.
\] The actual ring \(A\subset A_0\) has \(\operatorname{length}(A_0/A)=2\), since \(A_0\) has delta invariant two. Each branch projection of \(A\) is onto.

Let \(W_1\subset F^1/F^2=\mathbb C^3\) be its first-jet image. Each coordinate projection is nonzero. If \(W_1=\mathbb C^3\), Nakayama first gives \(\mathfrak m_AA_0=F^1\) and then, on the finite \(A\)-module \(A_0/A\), gives \(A=A_0\). This contradicts colength two. Thus \(\dim W_1=1\) or \(2\).

\par\medskip\noindent\textbf{2a. One-dimensional first jets}\par\label{a.-one-dimensional-first-jets}

Scale the branch parameters so \(W_1=\mathbb C(1,1,1)\). The two missing first jets use the entire colength, hence \(F^2\subset A\). The exact algebra is \[
 A=\mathbb C_{\rm diag}+\mathbb C(t_1,t_2,t_3)+F^2.
\] It is generated by \[
 x=(t_1,t_2,t_3),\quad y=(t_1^2,0,0),\quad z=(0,t_2^2,0).
\] Indeed \(y,z,x^2-y-z\) give all three quadratic branch ideals, and multiplication by \(x\) gives every higher order. The embedded model is exactly \[
       P=\bigl(y(y-x^2),yz,z(z-x^2)\bigr).             \tag{C}
\] The branches are \((t,t^2,0),(t,0,t^2),(t,0,0)\). Its differentiated generators give \(E_i=t_i^2K_i\) on all three branches. For example, on \((t,t^2,0)\) use the primitive basis \(dy-2t\,dx,dz\). Formula (A) gives \[
                         D\ge4+3\cdot3=13.            \tag{D}
\]

\par\medskip\noindent\textbf{2b. Two-dimensional first jets}\par\label{b.-two-dimensional-first-jets}

The three branch evaluations are nonzero linear forms on \(W_1\). If no two are proportional, their squares are independent: choose the first two as \(s,t\); the third is \(as+bt\), \(ab\ne0\), whose square supplies the nonzero \(st\) coefficient. Thus products give the whole second-jet space. Multiplication by one first jet nonzero on all branches then gives every later jet. This would leave only one missing positive-jet dimension, contrary to colength two.

Two evaluations are therefore proportional. Scale and label the parameters to obtain \(W_1=\{(a,a,b)\}\). Its products give the same two-dimensional second-jet space. The one remaining missing dimension occurs there: if the second-jet image were full, multiplication by a nonvanishing first jet would make every later one full. Hence \(F^3\subset A\), and the pure second-jet image is exactly \(\{(c,c,d)\}\).

A lift of \(W_1\) has quadratic difference \(\lambda a+\mu b\). The analytic changes \(t_i\mapsto t_i+c_it_i^2\), on the first two branches, remove \(\lambda a\). Thus the exact ring has common constants, with first and second coefficients satisfying \[
                     a_1=a_2,\qquad c_1-c_2=\mu a_3, \tag{E}
\] and arbitrary independent terms of order at least three.

For \(\mu\ne0\), the two elements \[
 x=(t_1,t_2,0),\qquad y=(\mu t_1^2,0,t_3)
\] generate this whole ring. The exact identities \[
 \begin{aligned}
 xy/\mu&=(t_1^3,0,0),\\
 x^3-xy/\mu&=(0,t_2^3,0),\\
 y^3-\mu(xy)^2&=(0,0,t_3^3),\\
 y^2-\mu x^2y&=(0,0,t_3^2)
 \end{aligned}
\] give the independent cubic ideals and every permitted quadratic jet. Multiplication by \(x\) or \(y\) gives all the higher branch series. This ring is planar, contrary to the present hypothesis.

Thus \(\mu=0\). The ring is generated by \((t_1,t_2,0),(0,0,t_3),(t_1^3,0,0)\). Naming the independent transverse coordinate \(z\) gives the exact embedded model \[
                       P=(xz,yz,y(y-x^3)).            \tag{F}
\] Its branches are \((t,0,0),(t,t^3,0),(0,0,t)\), and its conormal elementary divisors are \[
                         (1,3),\quad(1,3),\quad(1,1).
\] On the second branch use \(dy-3t^2dx,dz\); the gradients give \(t^3(dy-3t^2dx)\) and \(t\,dz\). Without restricting a symbolic-square coefficient of \(f\), (A) gives \[
                  D\ge4+\frac52+\frac52+\frac32
                       =\frac{21}{2}>9.               \tag{G}
\]

These are actual analytic ring isomorphisms. In each nonplanar case the three displayed generators form a basis of the curve cotangent space. Lift them to the ambient local ring. Their linear Jacobian is invertible, so the inverse-function theorem gives an ambient coordinate isomorphism to the stated model.

\par\medskip\noindent\textbf{3. Multiplicity four: one cusp and two smooth branches}\par\label{multiplicity-four-one-cusp-and-two-smooth-branches}

The singular branch has multiplicity two and delta at least one. Joining three branches costs at least two further delta units. If the branch delta is two, all remaining gluing is minimal, and the local ring is \[
 \mathbb C_{\rm diag}
 +\mathfrak m_{\rm singular}
 +t_2\mathbb C\{t_2\}+t_3\mathbb C\{t_3\}.
\] Its embedding dimension is at least \(2+1+1=4\), impossible in a smooth threefold. A larger branch delta is impossible. Thus this branch is the ordinary cusp \(\mathbb C\{t^2,t^3\}\).

Set \[
 A_0=\mathbb C_{\rm diag}
 +\mathfrak m_{\mathbb C\{t^2,t^3\}}
 +t_2\mathbb C\{t_2\}+t_3\mathbb C\{t_3\}.
\] Now \(A_0\) has delta three, and \(A\subset A_0\) has colength one. The one-dimensional \(A\)-module \(A_0/A\) is killed by \(\mathfrak m_A\), so \(\mathfrak m_AA_0\subset A\). Its elements have zero constant term, giving \(\mathfrak m_AA_0=\mathfrak m_A\). Hence \(\mathfrak m_A\) is an ideal of \(A_0\). The local algebra \(A_0/\mathfrak m_A\) has length two, so its maximal ideal squares to zero. Therefore \[
                         \mathfrak m_0^2\subset A.     \tag{H}
\]

The four cotangent coordinates of \(A_0\), denoted \(a,b,c,d\), are represented by the cusp functions \(t^2,t^3\) and the two smooth parameters. The three-dimensional subspace \(W=\mathfrak m_A/\mathfrak m_0^2\) is a hyperplane \[
                 \lambda a+\mu b+\nu c+\xi d=0.       \tag{I}
\] Its projection to \(a,b\) is onto, since the actual branch image is the full cusp ring, and the two smooth projections are nonzero. Ambient embedding dimension three gives \(\mathfrak m_A^2=\mathfrak m_0^2\). Thus products of \(W\) must generate \(\mathfrak m_0^2/\mathfrak m_0^3\), whose coordinates are \[
                              a^2,\quad ab,\quad c^2,\quad d^2.
                                                               \tag{J}
\]

\par\medskip\noindent\textbf{3a. A nonzero coefficient of the cusp cubic}\par\label{a.-a-nonzero-coefficient-of-the-cusp-cubic}

For \(\mu\ne0\), solve (I) for \(b\). Surjectivity of the cusp projection forces \((\nu,\xi)\ne(0,0)\). The three minimal ambient coordinate functions have parametrizations \[
 \begin{array}{c|c}
 \text{cusp}&(t^2-\lambda't^3,-\nu't^3,-\xi't^3)\\
 \text{second branch}&(0,t,0)\\
 \text{third branch}&(0,0,t).
 \end{array}
\] Their tangent directions are independent.

By P4R.15, \(f\in\mathfrak m^3\). Let \(G_3\) be its homogeneous cubic part. Since \(df\) vanishes on all three branches, \(dG_3\) vanishes on each independent tangent axis. Every cubic monomial other than \(xyz\) has a derivative nonzero on an axis. Hence \(G_3=cxyz\). On the \(x\)-axis its normal Hessian is \[
                         \begin{pmatrix}0&c\\c&0\end{pmatrix}.
\] A nonzero \(c\) gives a nonzero leading minor on the actual branch, contradicting generic rank one. Thus \(c=0\), and \(f\in\mathfrak m^4\).

The Hessian entries have order at least \(2m_i\) on branch \(i\), and \(b_i\ge m_i\). Its defect contribution is at least \(m_i+1\). Therefore \[
                            D\ge4+(4+3)=11>9.         \tag{K}
\]

\par\medskip\noindent\textbf{3b. A zero coefficient of the cusp cubic}\par\label{b.-a-zero-coefficient-of-the-cusp-cubic}

For \(\mu=0\), surjectivity in (J) forces \(\lambda,\nu,\xi\) all to be nonzero. Otherwise two of \(a,c,d\) are proportional on \(W\), so their squares are dependent, and (J) has rank at most three. That creates an extra minimal generator, contrary to embedding dimension three. Scale the branch parameters so that \(a+c+d=0\).

The pure hyperplane generators belong to \(A\) exactly by (H). They give the parametrizations \[
                       (t^2,t^2,t^3),\quad(t,0,0),\quad(0,t,0).
\] Their defining ideal is \[
                 P=(z(x-y),xy(x-y),z^2-x^2y).        \tag{L}
\] A Hilbert--Burch matrix is \[
 \begin{pmatrix}
 xy&z\\
 -z&-x\\
 0&-(x-y)
 \end{pmatrix}.
\] Its maximal minors give (L), with alternating signs. The ideal has codimension two and is Cohen--Macaulay. Its zero set is exactly the cusp and the two axes, and it is generically reduced on each branch. A nonzero nilpotent quotient supported at the origin would violate depth one, so it is the actual reduced intersection ideal.

On the cusp use the primitive conormal basis \[
                     k_1=dx-dy,\qquad k_2=3t\,dx-2dz.
\] The differentiated generators are \(t^3k_1,t^4k_1,t^3(tk_1-k_2)\). Thus its elementary divisors are \((3,3)\). On each smooth axis they are \((1,2)\). Here \(r=1\) on the cusp and zero on the axes. By (A), \[
                  D\ge4+\left(\frac92-1\right)+2+2
                       =\frac{23}{2}>9.               \tag{M}
\]

This exhausts multiplicity four. Together with Sections 1--2 it proves the strict local bound for every nonplanar three-branch curve with delta four. Only actual coordinate isomorphisms, conormal maps, and the rank-one Hessian were used; no boundary deformation has been introduced.

\par\medskip\noindent\textbf{4. Every nonordinary triple exceeds the budget}\par\label{every-nonordinary-triple-exceeds-the-budget}

For a planar three-branch point, P4.48 gives defect at least ten, including the ordinary plane triple point. At a nonplanar point with delta at least five, (B) and \(m\ge3\) give \[
                  D\ge5+3+\frac32=\frac{19}{2}>9.
\] Delta four was exhausted above.

At delta three the excess \(\delta-(r-1)\) is one. The actual analytic ring calculation in P4.40 gives precisely an ordinary plane triple, a plane tacnode with an independent transverse line, or a cusp with two independent lines. The first is excluded by P4.48. The second has defect at least eleven by P4R.16. The third has embedding dimension four and cannot occur in the ambient smooth threefold.

Finally a three-branch point with delta two has zero excess. Its ring equals \(\mathbb C_{\rm diag}+F^1\), generated by the independent branch parameters. It is exactly the three coordinate axes, with ideal \((xy,xz,yz)\). Thus every nonordinary three-branch point is excluded from the total defect budget nine. The only remaining local possibility is the ordinary nonplanar triple treated by P4.48 and the separate full-quartic conductor calculation.

\end{proof}

The analytic model identities, primitive conormal lattices,
Hilbert--Burch minors, cubic obstruction, and jet generators are
checked exactly in \path{code/part4/a3_delta4.py}.


\subsection{The horizontal image in the unrestricted ruled remainder}

\begin{statement}{P4R.18}{Proved}
In the remaining positive-Euler ruled case of Part~4.2, a
noncontracted horizontal pole section maps with degree two or
three to a rational normalized conductor component whose generic
transverse germ is $D_4$. This assertion imposes no transverse
type or conductor-exponent restriction on the other image component.
\end{statement}
\begin{proof}
Write $H$ for the horizontal section, $g=g(H)\ge1$, and
$C_H$ for its image in the reduced conductor $C$. The preceding
ruled reduction gives two rational normalized components of $C$
and one point-identification cycle. The nonconstant map
\[
                         H\longrightarrow C_H^\nu=\mathbb P^1
\]
has degree at least two: a degree-one finite map between smooth
curves is an isomorphism, whereas their genera differ. Each
generic point of $H$ has conductor exponent two, by the actual
pole divisor of P4R.6. Thus the generic transverse plane curve
along $C_H$ has at least two normalization branches of exponent two.

For its branches $B_i$ the exact conductor formula is
\[
             c_i=2\delta(B_i)+\sum_{j\ne i}I(B_i,B_j).
\]
A singular branch with $c_i=2$ permits no other branch,
so the two specified branches are smooth. With two total
branches their intersection multiplicity is two, giving $A_3$.
With three total branches, the first two exponents force all
three pairwise intersection multiplicities to equal one;
the third branch is then smooth as well, giving $D_4$.
Four branches would give exponent at least three on every
smooth branch. These possibilities are exhaustive. In type
$A_3$, $H$ has degree exactly two and supplies both generic
normalization preimages. In type $D_4$ its degree is two or
three. We exclude the former transverse type next.

Assume that the transverse germ along $C_H$ is $A_3$.
Let $C_{H,\rm sch}$ be the primary component of the actual
image conductor scheme supported on $C_H$. The primary
restriction construction of P4R.13 applies before its
elliptic-genus calculation, without change. More explicitly,
P4.33 gives rational singularities of the normalization when
$H$ is not contracted, so finite duality and the canonical
trace give the injections
\[
 \mathcal O_{C_{\rm sch}}\longrightarrow
 \nu_*(\mathcal O_N/I_N)\longrightarrow\rho_*\mathcal O_D.
\]
The exact blowup relation $D'=\beta^*D''-E$ identifies
global functions on $D$ with those on $2H_0+V_0$ of the
minimal ruled model. The kernel of restriction to $2H_0$
has successive vertical-fibre layers
$\mathcal O_{\mathbb P^1}(-2)$, with no global sections.
Hence a global function is determined at the generic doubled
horizontal section. It follows that restriction from
$C_{\rm sch}$ to $C_{H,\rm sch}$ is injective on $H^0$
and surjective on $H^1$: the first assertion follows from
those injections and generic determination, and the second
from the dimension-one kernel of the restriction sheaf map.
Therefore
\[
                 g-1=\chi(\mathcal O_{C_{\rm sch}})
                    \le\chi(\mathcal O_{C_{H,\rm sch}}).
                                                            \tag{P4R.18.1}
\]
Every arrow here is the primary restriction of the same
conductor algebra; no other conductor component is deleted
from the left side.

We evaluate the right side on $C_H^\nu=\mathbb P^1$.
Use the actual differentiated conormal, the saturated Hessian
line and the rank-one nilradical of P4R.10. The kernel $K$
of the normalized derivative has degree $2+\deg R$, since
$c_1(Y)\cdot C_H=0$. The Hessian line has degree
$\tfrac12\sum d_p$. Riemann--Roch and the normalization
sequence give exactly
\[
 \chi(\mathcal O_{C_{H,\rm sch}})
 =4-\sum_p\left[\delta(C_H,p)
      +\sum_{i\mapsto p}(d_i/2+b_i-r_i)+q_p\right].
                                                            \tag{P4R.18.2}
\]
All terms are those of this primary curve and its primary
conductor. They remain defined at its intersections with
the other component, because the pure conductor quotient
extends its generic containment of the symbolic square.

Every singular point of $C_H$ contributes at least three
to the sum in (P4R.18.2). At a planar point the ideal is
$(u,h)$, its symbolic square equals its ordinary square,
and the actual equation belongs to that square. The
rank-one symmetric-square estimate of P4.42 and the exact
plane-curve different calculation give a contribution at
least $3\delta(C_H,p)+q_p\ge3$.
At a nonplanar point every defining equation has order at
least two. On a branch of multiplicity $m_i$ its projected
conormal order is at least $m_i$. The point is singular
on the entire critical curve, so P4R.15 makes the ambient
Hessian zero and gives $d_i\ge m_i$. Since $r_i=m_i-1$,
the contribution is at least
\[
                      \delta(C_H,p)+\sum_i(1+m_i/2),
\]
which is at least three at a singular curve point. The
same arguments give nonnegative contributions everywhere;
at a smooth primary point its conormal is the full kernel
and its only possible contribution is the nonnegative
Hessian zero order divided by two.
Combining (P4R.18.1)--(P4R.18.2) gives total primary
defect at most $5-g\le4$. Thus $C_H$ has at most one
singular point, and two different singular points are
impossible even in the presence of arbitrary transverse
types on the other global component.

We count the possible branch values of the degree-two map
$H\to\mathbb P^1$ with the identifications retained.
The normalization graph formula is
\[
                         \sum_{p\in C}(r_p-1)=2.
\]
Connectedness leaves either two two-branch points or one
three-branch point. In the two-point case, two intersections
of the global components give two points on $C_H^\nu$.
There can be at most one additional unibranch singularity
of $C_H$, so these supply at most three possible branch
values. The other connected possibility is a self-identification
on one component and one intersection of the two components.
A self-identification on $C_H$ supplies two normalization
points and already costs at least three in the primary
defect. It permits no additional singular point, so with
the intersection there are again at most three possible
branch values. A self-identification on the other component
supplies no additional point on $C_H^\nu$.

At a three-branch point, the two global components must
supply two and one branches. Thus $C_H^\nu$ has at most
two points over it. Two points already make $C_H$ singular
there and exclude any further singular point by the defect
bound. One point permits at most one further unibranch
singularity. This case also has at most three possible
branch values.

There is no ramification elsewhere. At a point where the
entire critical curve is smooth, P4R.14 gives precisely
the product tacnode or $X(Y^2+XZ)$. Both have two distinct
normalization points over that point of the conductor.
Their two generic preimages lie on $H$, so properness
places both special preimages on $H$ as well. The
degree-two map is therefore unramified there. We have
accounted for every possible ramification value.
Riemann--Hurwitz for a degree-two map from a genus-$g$
curve requires $2g+2\ge4$ distinct branch values; each
ramified fibre consists of its single point of index two.
This contradicts the count of at most three. Thus the
generic germ along $C_H$ must be $D_4$, with degree two
or three as proved at the start.
\end{proof}

The exact connected identification patterns and their primary
defect/branch-value counts are enumerated in
\texttt{code/part4/horizontal\_primary.py}.


\subsection{A triple spectral root with nonzero cubic tensor}

\begin{statement}{P4R.19}{Proved}
Let $f$ be the actual reduced local equation of the remaining
positive ruled support, and let $C_0$ be a smooth local branch
of its critical curve with generic transverse type $D_4$.
At a point where the cubic normal tensor along $C_0$ is
nonzero, that tensor is not the cube of a linear form.
Other critical-curve branches at the point are retained,
without a restriction on their generic transverse types.
\end{statement}
\begin{proof}
Choose coordinates $(x,y,z)$ with $C_0=(x=y=0)$.
Generic normal multiplicity three gives
$f\in(x,y)^{(3)}$. Powers of this smooth
complete-intersection ideal are primary, so
$f\in(x,y)^3$ at the specified point.

Suppose its nonzero cubic normal tensor is a cube.
A linear normal-coordinate change makes it $x^3$, up
to a nonzero scalar. Weierstrass preparation in $x$
gives the exact equality
\[
 f=U(x,y,z)\bigl[x^3+a_2(y,z)x^2+a_1(y,z)x+a_0(y,z)\bigr],
\]
where $U$ is a holomorphic unit. The bracket belongs to
$(x,y)^3$, so $a_2,a_1,a_0$ are divisible by $y,y^2,y^3$,
respectively. Depressing the monic cubic is a holomorphic
coordinate change preserving $(x,y)$ and yields
\[
 f=U(x,y,z)\bigl[x^3+y^2B(y,z)x+y^3C(y,z)\bigr],
 \qquad B(0,0)=C(0,0)=0.                         \tag{R19.1}
\]
The unit remains in this equality and in the scaling
comparison below.

Set
\[
 k=\operatorname{ord}_z B(0,z),\qquad
 l=\operatorname{ord}_z C(0,z),\qquad
 \rho=\min(k/2,l/3),
\]
where the order of an identically zero series is infinity.
At least one order is finite: otherwise all generic normal
cubics would be $x^3$, contrary to generic $D_4$.
Thus $0<\rho<\infty$. Choose $L>3\rho$ and positive
rational weights
\[
       \operatorname{wt}(x)=L+\rho,\qquad
       \operatorname{wt}(y)=L,\qquad
       \operatorname{wt}(z)=1.
\]
Multiplication by a common denominator gives positive
integer weights. The weighted degree of the first face
is $3L+3\rho$. Terms of $B$ containing $y$ have weight
gap at least $L-2\rho>0$, and those of $C$ containing
$y$ have gap at least $L-3\rho>0$. Consequently the
leading polynomial, apart from the nonzero factor $U(0)$,
is exactly one of
\[
\begin{array}{ll}
 P=x(x^2+a y^2z^k),\quad a\ne0,
       & k/2<l/3,\\[2pt]
 P=x^3+b y^3z^l,\quad b\ne0,
       & l/3<k/2,\\[2pt]
 P=Q(x,yz^r),\quad
 Q(X,V)=X^3+aXV^2+bV^3,\quad a,b\ne0,
       & k=2r, l=3r, r\ge1.
\end{array}                                                   \tag{R19.2}
\]
In the last case $r$ is an integer. The depressed cubic
$Q$ is not a cube: a cube with $X^3$ coefficient one and
$X^2V$ coefficient zero would have $a=b=0$. A double
factor of $Q$ is allowed throughout the argument.

\par\medskip\noindent\textbf{The compact class used in the comparison.}
Let $B$ be a connected smooth complex curve and let
$h:B\to\mathbb C$ have two distinct simple zeros $p,q$.
Join them by a smooth embedded arc $\gamma$ whose
interior contains no zero. On the smooth surface
\[
             Z_h=\{(b,u,v)\in B\times\mathbb C^2:
                                      uv=h(b)\},
\]
consider the locus over $\gamma$ with $|u|=|v|$.
It is a circle over each interior point and collapses
to a point at each end. Near an endpoint choose the
arc so that $h$ is a fixed nonzero complex multiple
of a positive real parameter $s$. Then
$u=\sqrt{s}\,e^{\mathrm i\phi}$ and
$v=c\sqrt{s}\,e^{-\mathrm i\phi}$ are the smooth
polar coordinates of a real disc in the local $uv$
chart. The locus is thus a smoothly embedded sphere.

It is Lagrangian for a product K\"ahler form.
Writing $u=R e^{\mathrm i\phi}$,
$v=R e^{\mathrm i(\psi-\phi)}$, the sum of the two
fibre area forms restricts to
$d(R^2/2)\wedge d\psi=0$, since $R,\psi$ depend
only on the arc parameter. The base area form also
restricts to zero. Multiplication by $\mathrm i$
identifies tangent and normal bundles, with the
opposite normal orientation relative to the complex
ambient orientation. Its self-intersection is
therefore $-\chi(S^2)=-2$.
For the compactly supported Poincar\'e dual $\alpha$,
\[
                         \int_{Z_h}\alpha^2=-2.             \tag{R19.3}
\]

If $\pi:W\to Z_h$ is a finite proper holomorphic map
of degree $d>0$ between smooth complex surfaces, then
$\pi^*\alpha$ is compactly supported and
\[
              \int_W(\pi^*\alpha)^2
                       =d\int_{Z_h}\alpha^2=-2d.            \tag{R19.4}
\]
The equality follows from the proper oriented degree;
off the analytic branch locus the map has $d$ positive
sheets, and the branch locus has measure zero for
integration of compactly supported top forms.
Poincar\'e duality gives a compact homology two-class
with that square. Its representative need not itself
be a sphere.

\par\medskip\noindent\textbf{The exact finite maps.}
In the first case of (R19.2), $P=1$ forces $x\ne0$
and gives
\[
           y^2z^k=h(x),\qquad
           h(x)=\frac{1-x^3}{ax},\qquad x\in\mathbb C^*.
\]
The map
\[
 (x,y,z)\longmapsto(x,u=y^2,v=z^k)
\]
is finite proper of degree $2k$ onto $uv=h(x)$.
The target is smooth: at $u=v=0$ the zeros of $h$
are simple, and elsewhere a derivative in $u$ or
$v$ is nonzero. The source is smooth by the weighted
Euler identity on $P=1$. Join two of the three cube
roots of one by an arc in $\mathbb C^*$ avoiding the
third. Equations (R19.3)--(R19.4) give a compact
two-class of square $-4k$ on the leading fibre.

In the second case the exact finite proper map is
\[
 \begin{split}
 \{x^3+b y^3z^l=1\}&\longrightarrow
                 \{uv=(1-x^3)/b\},\\
 (x,y,z)&\longmapsto(x,u=y^3,v=z^l).
 \end{split}
\]
It has degree $3l$, and the base is $\mathbb C$.
The same three simple zeros supply a pulled-back
compact class of square $-6l$.

For the balanced case let
\[
                 B_Q=\{(X,V)\in\mathbb C^2:Q(X,V)=1\}.
\]
It is smooth by the ordinary Euler identity and is
connected even when $Q$ has a double factor. Its
projection to $\mathbb P^1$ minus the distinct zero
directions of $Q$ is an unramified cyclic cover of
degree three. Around a zero direction of multiplicity
one or two its monodromy generates $\mathbb Z/3$.
Such a direction exists because $Q$ is not a cube,
so the cover is connected.

The function $V$ on $B_Q$ has precisely the three
zeros $(\zeta,0)$ with $\zeta^3=1$, and they are
simple because $Q_X(\zeta,0)=3\zeta^2\ne0$.
The surface
\[
             Z_Q=\{(b,u,w):b\in B_Q,\ uw=V(b)\}
\]
is smooth. The exact map
\[
 \begin{split}
       \{Q(x,yz^r)=1\}&\longrightarrow Z_Q,\\
       (x,y,z)&\longmapsto
                    \bigl(b=(x,yz^r),u=y,w=z^r\bigr)
 \end{split}                                                   \tag{R19.5}
\]
is finite proper of degree $r$: $z$ satisfies the
monic equation $z^r=w$ and the other coordinates are
retained. An arc in $B_Q$ joining two $V$-zeros and
avoiding the third supplies a compact class of square
$-2r$ on the source. In particular this argument
retains $Q=(X-\lambda V)^2(X+2\lambda V)$,
$\lambda\ne0$, without a perturbation of its
repeated factor.

\par\medskip\noindent\textbf{Transport to the actual local fibre.}
Let $(w_x,w_y,w_z)$ be the positive integer weights
and let $d$ be the resulting degree. Put
\[
 D_\epsilon(x,y,z)=
       (\epsilon^{w_x}x,\epsilon^{w_y}y,\epsilon^{w_z}z),
 \qquad
 G_\epsilon=\epsilon^{-d}f\circ D_\epsilon.
\]
These actual functions converge with derivatives on
compact sets to $U(0)P$. The weighted Euler identity
makes $dP$ nonzero on $P=1$, including the repeated-factor
case on its zero set. Choose compact representatives
of the two-class above and of the transverse
perturbation computing its square. On a compact
neighbourhood of these chains, the gradient-flow
comparison of P4.52 solves
\[
       dG_\epsilon(V_\epsilon)
                         =-\partial_\epsilon G_\epsilon.
\]
The denominator is bounded away from zero on a
slightly larger neighbourhood; a cutoff gives a
flow for sufficiently small positive $\epsilon$.
It transports the chains and their intersections
to $G_\epsilon=U(0)$. Scaling by $D_\epsilon$ then
embeds them in the original fixed Milnor ball on
\[
                         f=U(0)\epsilon^d.
\]
The oriented intersection number is preserved and
equals $-4k$, $-6l$, or $-2r$. Each is nonzero,
contradicting P4.51. This proves the assertion.
\end{proof}

\par\medskip\noindent\textbf{The two-direction exception is retained.}
The assertion cannot be enlarged to every spectral
collision. The reduced equation
\[
                       f=y(x^2+y^2z^2)
\]
has generic $D_4$ along $x=y=0$ and generic $A_1$
along $x=z=0$. Only two spectral directions collide
at their intersection. Its non-plane normalization
sheets $x=\pm\mathrm i yz$ have pole divisor
$2(y=0)+(z=0)$, with simple normal crossings.
Its affine fibre $f=1$ is exactly $(\mathbb C^*)^2$:
\[
 \begin{split}
       a&=x+\mathrm i yz,\qquad b=x-\mathrm i yz,\\
       x&=(a+b)/2,\qquad y=(ab)^{-1},\qquad
       z=(a-b)ab/(2\mathrm i).
 \end{split}
\]
The displayed maps are inverse and substitution gives
$f=1$. Its intersection form is zero: the product
of two unit circles generates its second homology,
and a radial translation gives a disjoint homologous
copy. More generally the fibre of $y(x^2+y^2z^n)$
becomes $u^2+z^n=t^3$, $t\ne0$, through
$t=1/y$, $u=x/y$. Positive radial dilation gives
a product real direction, along which any compact
cycle can be displaced from another compact cycle.
Its intersection form is also zero. Thus the
other generic transverse type remains unrestricted.

The exact weighted gaps, all three finite-map
equations, their degree-square identities, and the
inverse maps of the two-direction exception are
executed in \path{code/part4/d4_spectral_triple.py};
its \path{.json} and \path{.out} outputs report PASS.


\subsection{The three-branch image point with a double horizontal cover}

\begin{statement}{P4R.20}{Proved}
In the remaining degree-two horizontal sector after P4.55,
suppose that the generic transverse curve along the other
image component $B$ has at least three normalization branches.
The unique image-graph cycle cannot be concentrated at one
point with three local critical-curve branches, namely the
smooth branch $A$ and two local branches of $B$.
\end{statement}
\begin{proof}
Write these two local branches as $B_+,B_-$ and the
point as $p$. Let $W_1,\ldots,W_c$ be the normalized
conductor components mapping onto $B$, with degrees
$d_i>0$ over $B^\nu$. Put
\[
                         r_B=\sum_i d_i\ge3.
\]
This is the number of generic normalization branches
of the full transverse plane curve along $B$.
The normalization conductor has components
$H,R,W_1,\ldots,W_c$ and a tree graph. Its total
point-identification contribution is $c+1$.
The horizontal section $H$ maps with degree two to
$A$, while the rational vertical curve $R$ maps
with degree one. Distinct points of $H$ are not
identified in the normalization $N$. Distinct
points of $R$ are not identified either: a
contracted vertical tree effecting such an
identification would create a cycle in the ruling
fibre. All conductor branches below are counted
on their normalization.

We first record the exact spectral bound used in
the proof. Let $E=N_{A/Y}^*$, so $\deg E=2$ by
adjunction and $c_1(Y)\cdot A=0$. On $\mathbb P(E)$
use the quotient convention, tautological class
$\xi$, and fibre class $F$. Thus
\[
          \xi^2=2,\qquad \xi F=1,\qquad F^2=0,
          \qquad K_{\mathbb P(E)}=-2\xi.
\]
The actual normal derivative maps of P4.55 identify
$H$ with the normalization of the degree-two
spectral component $\Gamma_H$, and identify the
third sheet with a section $\Gamma_R\cong A$.
Remove the common zero divisor of the cubic tensor,
of degree $b\ge0$. For an integer $t$ write
\[
       \Gamma_R=\xi-tF,\qquad
       \Gamma_H=2\xi+(t-b)F.
\]
If $\delta_H$ is the normalization colength of
$\Gamma_H$ and $n=\Gamma_H\cdot\Gamma_R\ge0$,
adjunction and the intersection product give
\[
 \begin{split}
       g(H)+\delta_H&=p_a(\Gamma_H)=1+t-b,\\
       n&=4-t-b,\\
       g(H)+\delta_H+n&=5-2b.
 \end{split}                                                   \tag{R20.1}
\]
These are the actual residual components of the
cubic, with their common zero divisor retained.
Since $g(H)\ge1$, they imply $b\le2$.

\par\medskip\noindent\textbf{The exact graph when $c\ge2$.}
Let $r_A$ be the number of normalized conductor
points over the branch $A$ at $p$. It is two if
$H\to A$ ramifies there, and three otherwise.
Let $r_+,r_-$ count the points over $B_+,B_-$.
Every $W_i$ is surjective to $B^\nu$, so
$r_+,r_-\ge c$. Let $n_p$ be the number of
normalization points of $S$ over $p$.
Their nonempty supports in $\{A,+,-\}$ pairwise
intersect: distinct local hypersurface components
intersect in a curve.

A singleton support $\{A\}$ forces every point to
carry $A$. Then $n_p\le r_A$ and the graph
contribution is at least $r_++r_-\ge2c>c+1$.
A singleton $\{+\}$ forces every point to carry
$+$ and gives contribution at least
$r_A+r_-\ge c+2$; the same applies to $\{-\}$.
Thus every support has at least two elements, and
\[
 \text{contribution at }p
       \ge\frac{r_A+r_++r_-}{2}\ge c+1.
\]
The whole global budget is $c+1$. Equality is
forced throughout:
\[
                 r_A=2,\qquad r_+=r_-=c,\qquad n_p=c+1.
                                                               \tag{R20.2}
\]
Every point over $p$ has exactly two normalized
conductor branches, with different image labels,
and all graph identifications occur at $p$.
The horizontal cover ramifies over $p$.
Each $W_i$ has precisely one point over each of
$B_+,B_-$ and has valence two; $H$ and $R$ have
valence one. The connected tree is therefore
the path
\[
                    H-W_1-W_2-\cdots-W_c-R.             \tag{R20.3}
\]
Choose the label $+$ at the $H$ endpoint. The next
point on $W_1$ has label $-$. Every internal point
has support $\{+,-\}$, so it joins that point to
the $+$ point of $W_2$. Induction gives the $+$
point on the left of every $W_i$ and its $-$
point on the right. The $R$ endpoint carries $-$.

The endpoint hypersurface factors, denoted
$f_H,f_R$, thus have supports $\{A,+\}$ and
$\{A,-\}$. The $c-1$ internal factors $f_j$
have supports $\{+,-\}$. These are all the local
irreducible hypersurface factors.

\par\medskip\noindent\textbf{The actual length-two ribbon.}
The set-theoretic intersection of $f_H$ and $f_R$
is $A$ alone. Along generic $A$, the first factor
supplies two ordinary-triple branches and the
second supplies the third, so their intersection
multiplicity is two. Their one-dimensional
complete-intersection ring is Cohen--Macaulay,
without a zero-dimensional embedded component.
A coordinate $z$ on the smooth $A$ is a
nonzerodivisor. The ring is finite free of rank
two over $\mathbb C\{z\}$, and its quotient by
$z$ has length two.

If both factors had no linear term, their
restrictions modulo $z$ would belong to
$(x,y)^2$. The Artinian quotient would contain
the independent classes $1,x,y$, contrary to
length two. At least one factor is smooth.
It cannot be $f_H$: a smooth surface is its own
normalization and supplies only one generic
preimage of $A$, whereas this factor supplies
the ramified double cover. Thus $f_R$ is smooth.
Take the actual ambient coordinate $u=f_R$.
On $u=0$, the divisor of $f_H$ is exactly twice
the smooth $A$, so in coordinates
$A=(u=v=0)$ it has restriction $v^2$ times a
unit. Weierstrass preparation in $v$, followed
by depression and using $f_H\in(u,v)^2$, gives
the exact expression
\[
                    f_H=\text{unit}\,(v^2+u^2g(u,z)).     \tag{R20.4}
\]
All changes preserve $A$. The full cubic normal
tensor is therefore, up to a nonvanishing scalar,
\[
        \left(\prod_j f_j(0,0,z)\right)
                     u\bigl(v^2+u^2g(0,z)\bigr).          \tag{R20.5}
\]
The coefficient of $uv^2$ in the last factor is
one, so the common zero order is exactly the
sum of the orders of $f_j|_A$.

\par\medskip\noindent\textbf{The pole on the smooth rational branch.}
On the smooth branch $f_R=0$, its conductor is
the restriction of the product of the other
factors. The factor $f_H$ contributes $2A$.
Each $f_j$ cuts out only $B_-$ there. This
normalization branch is already smooth, and
the pole support on its minimal resolution is
SNC. Hence $A$ and its $B_-$ curve are smooth
and transverse. Choose the latter's equation
$z=0$ on this surface, with $z$ still a
coordinate on $A$, and extend it to the ambient
threefold. Then
\[
             f_j|_{f_R=0}=\text{unit}\,z^{e_j},
                  \qquad e_j\ge1.
\]
Their sum $e=\sum e_j$ is the generic conductor
exponent of the global component $W_c$ on
this branch. By (R20.5) it is also the local
common cubic-zero order. The branch on
$f_R=0$ is smooth; the plane-curve conductor
formula therefore gives
\[
          e=\sum_{\text{other branches}}I
                         \ge r_B-1\ge2.
\]
Combining $b\ge e$ with (R20.1) forces
\[
     b=e=2,quad r_B=3,quad
     g(H)=1,quad \delta_H=n=0.                         \tag{R20.6}
\]
Thus $\Gamma_H$ is smooth and disjoint from
$\Gamma_R$, and there is no further common
zero of the cubic.

\par\medskip\noindent\textbf{The exact normalization at the ramified endpoint.}
The cover $H\to A$ ramifies at $p$ by (R20.2).
Equation (R20.4) presents its spectral cover as
$w^2+g(0,z)=0$. After translating its double
root, smoothness from (R20.6) gives
$g_z(0,0)\ne0$. Hence
\[
 \begin{split}
        N_H&=\{w^2+g(u,z)=0\},\\
        (u,w,z)&\longmapsto(u,v=uw,z)
 \end{split}                                                   \tag{R20.7}
\]
is the actual normalization near that point.
It is smooth, finite, and birational onto the
branch (R20.4); off $u=0$ its inverse is
$w=v/u$. Coordinates $(u,w)$ solve
$z=\psi(u,w)$, with $z|_H$ of order two in
$w$. The minimal normalization resolution
has no exceptional curve over this smooth
point, and $H$ is $u=0$.

The other conductor branch $W_1$ meets $H$
transversely by SNC. Thus $u$ is a coordinate
on $W_1$ and has order one on its image
$B_+$. By (R20.2), $W_1$ has only one point
over the normalized $B_+$ point, so its
global degree is $d_1=1$. The same smooth
SNC argument on $f_R=0$ gives $d_c=1$.
Since $\sum d_i=r_B=3$, the case $c=2$
would have sum two, while $c>3$ is impossible.
Therefore
\[
                 c=3,\qquad d_1=d_2=d_3=1.              \tag{R20.8}
\]
There are exactly two internal factors, and
$e_1=e_2=1$.

\par\medskip\noindent\textbf{The conductor mismatch.}
Each $f_j$ intersects $f_H$ only along $B_+$.
Its pullback to the actual smooth normalization
(R20.7) consequently has divisor supported
only on $W_1$. Its restriction to $H$ vanishes
to order two: the order on $A$ is one and
$H\to A$ has ramification index two.
Transversality of $H$ and $W_1$ makes the
divisor exactly $2W_1$. Thus the generic
intersection multiplicity of $f_H$ with
each internal factor along $B_+$ is two.

The path (R20.3) identifies the same global
component $W_3$ on the smooth $R$ branch
over $B_-$ and on an internal factor over
$B_+$. At $B_-$ its conductor exponent is
$e=2$. At $B_+$ it has intersection
multiplicity two with $f_H$ and positive
intersection multiplicity with the other
internal factor. Its conductor exponent
is at least three. These cannot be two
values of the coefficient on one global
irreducible conductor component.

The full marking is also explicit. Over
$B_+$ the factors $(f_H,f_1,f_2)$ carry
$(W_1,W_2,W_3)$; over $B_-$ the factors
$(f_1,f_2,f_R)$ carry the same ordered
global components. If $p_+,p_-\ge1$ are
the intersections of the internal factors,
the conductor vectors are
\[
         (4,2+p_+,2+p_+),\qquad
         (1+p_-,1+p_-,2).
\]
There is no missing intrinsic branch-delta
term: the exponent two of the smooth $R$
branch forces its two intersections to
equal one, which makes both other generic
$B$ branches smooth as well. The third
entries already give the contradiction.

\par\medskip\noindent\textbf{The case $c=1$.}
Now the conductor graph has $H,R,W$ and
budget two. The points of $W$ over $B_+$
and $B_-$ cannot be identified in $N$,
since that creates a loop in the conductor
tree. Points carrying the two labels are
distinct; their supports intersect, so
both carry $A$. Each contributes at least
one, exhausting the graph budget. Each
has exactly one $A$ and one $B$ branch,
and no other identification occurs.

If $H$ ramifies over $p$, these two
$A$ branches are $H$ and $R$. There are
only the factors with supports
$\{A,+\}$ and $\{A,-\}$. The length-two
ribbon argument makes $f_R$ smooth.
No other factor contains $B_-$, and this
smooth branch has no singularity there.
Thus $B_-$ cannot be a critical curve.

If $H$ is unramified, the three normalized
$A$ preimages are distinct. Two occur
at the $B$ attachments and the third
is an $A$-only normalization point.
There are three factors, each supplying
one generic $D_4$ branch. Every pair
meets only along $A$, with multiplicity
one. Its complete-intersection curve
is the reduced smooth $A$: it is pure,
generically reduced, and has no other
component. Modulo $z$, length one forces
independent normal linear terms.
All three factors are smooth near $p$
and their only pairwise intersection is
$A$. Neither $B_+$ nor $B_-$ can be a
critical branch. This final contradiction
completes the proof.
\end{proof}

The exact spectral intersections, length-two
ribbon basis, normalization equation, scalar
cubic coefficient, endpoint-degree list and
marked conductor vectors are executed in
\path{code/part4/degree_two_triple_point.py}.
Its \path{.json} and \path{.out} certificates
report PASS.

